% =====================================================================
%  The Irreducible Boundary of a Bounded Resolvable Space
%  A self-contained theory of individuation, boundary thickness,
%  and the resolution floor.
% =====================================================================
\documentclass[11pt,twocolumn,a4paper]{article}

\usepackage[utf8]{inputenc}
\usepackage[T1]{fontenc}
\usepackage{lmodern}
\usepackage[a4paper,margin=2.5cm]{geometry}
\usepackage{amsmath,amssymb,amsthm,mathtools}
\usepackage{bm}
\usepackage{mathrsfs}
\usepackage{booktabs}
\usepackage{array}
\usepackage{enumitem}
\usepackage{microtype}
\usepackage{graphicx}
\usepackage[round,sort&compress]{natbib}
\usepackage{tikz}
\usetikzlibrary{arrows.meta,positioning,calc,decorations.pathmorphing,fit}
\usepackage[colorlinks=true,linkcolor=blue!45!black,
            citecolor=blue!45!black,urlcolor=blue!45!black]{hyperref}
\usepackage{cleveref}

% ---- Theorem environments ----
\newtheoremstyle{thmstyle}{\topsep}{\topsep}{\itshape}{}{\bfseries}{.}{.5em}{}
\theoremstyle{thmstyle}
\newtheorem{theorem}{Theorem}[section]
\newtheorem{lemma}[theorem]{Lemma}
\newtheorem{proposition}[theorem]{Proposition}
\newtheorem{corollary}[theorem]{Corollary}
\theoremstyle{definition}
\newtheorem{definition}[theorem]{Definition}
\newtheorem{axiom}{Axiom}
\newtheorem{principle}[theorem]{Principle}
\newtheorem{example}[theorem]{Example}
\newtheorem{construction}[theorem]{Construction}
\theoremstyle{remark}
\newtheorem{remark}[theorem]{Remark}
\newtheorem{notation}[theorem]{Notation}

\crefname{axiom}{Axiom}{Axioms}
\Crefname{axiom}{Axiom}{Axioms}
\crefname{principle}{Principle}{Principles}
\Crefname{principle}{Principle}{Principles}

% ---- Shorthand ----
\newcommand{\R}{\mathbb{R}}
\newcommand{\N}{\mathbb{N}}
\newcommand{\Z}{\mathbb{Z}}
\newcommand{\Space}{\Omega}
\newcommand{\Bdy}{\partial}
\newcommand{\interior}{\mathrm{int}}
\newcommand{\cl}{\mathrm{cl}}
\newcommand{\diam}{\mathrm{diam}}
\newcommand{\dist}{\mathrm{dist}}
\newcommand{\Leb}{\mathcal{L}}
\newcommand{\Haus}{\mathcal{H}}
\newcommand{\thick}{\beta}
\newcommand{\reso}{\delta}
\newcommand{\floor}{\thick}
\newcommand{\Part}{\mathscr{P}}
\newcommand{\cell}{\mathsf{c}}
\newcommand{\eps}{\varepsilon}
\newcommand{\Neg}{\mathcal{N}}
\newcommand{\Ver}{\mathsf{V}}
\newcommand{\norm}[1]{\lVert#1\rVert}
\newcommand{\abs}[1]{\lvert#1\rvert}
\newcommand{\dd}{\,\mathrm{d}}

\title{\textbf{The Irreducible Boundary of a Bounded Resolvable Space:\\[2pt]
Individuation, Boundary Thickness, and the Resolution Floor}}

\author{%
Kundai Farai Sachikonye\\
\small Technical University of Munich\\
\small Chair of Experimental Bioinformatics\\
\small AIMe Registry for Artificial Intelligence\\
\small \texttt{kundai.sachikonye@wzw.tum.de}}

\date{\today}

\begin{document}
\maketitle

% =====================================================================
\begin{abstract}
\noindent
We study \emph{individuation} in a bounded resolvable space: the act of
separating a part from the whole when no element of the space is
distinguished in advance. We prove that, in a bounded space carrying a
finite resolution, individuation can proceed only by negation---a part is
fixed as the complement of everything it is not---and that every such
separation deposits content of strictly positive measure into a boundary.
This content, the \emph{resolution floor} $\thick>0$, is the central
invariant of the paper. We establish a \emph{Boundary--Thickness Theorem}
(no separator between a part and its complement has zero content), a
\emph{Non-Instantaneity Theorem} (a separation that yields a genuine part
cannot be completed in zero steps, and what the process leaves untouched is
an irreducible residue), and a \emph{Detectability Theorem} (a part is
distinguishable exactly when its own scale exceeds its boundary thickness;
sub-floor parts are engulfed by their separator). The classical sorites is a
corollary: there is no sharp count at which an aggregate becomes
distinguishable, because the separating boundary has positive width and
distinguishability is emergence from that width. We give three independent
derivations of the floor---representational, cardinality-theoretic, and a
cost argument by which compressing a separator to zero content diverges---and
prove they bound the same constant. We then separate two relations a region may
bear: being \emph{located} by an agent that individuates it, which is costly and
whose cost grows with the logarithm of the region's rarity, and being
\emph{visited} by a measure-preserving traversal, which is free and scales with
rarity only in frequency; from this asymmetry we classify three invalid
inferences that transfer a cost of locating onto being reached. We then identify
the deepest obstruction:
to \emph{verify} that a part is what it is requires deciding its membership
in the very collection that constitutes it (its negation), a self-referential
question with no consistent value. We prove, by a self-contained diagonal
argument, that no verifier can be both total and closed under its own diagonal;
\emph{under the explicit hypothesis} that an agent's verifying behaviour is itself
realisable as a region---which embedding makes natural but which we do not claim
to discharge in general---the irreducible residue is therefore not a locatable set
but a self-referential impossibility that cannot be exhibited, only bounded below
by $\thick$. The positivity of $\thick$ itself is established earlier and
unconditionally. From this we obtain a \emph{Non-Return Theorem}---the count of
committed separations is monotone non-decreasing under every traversal, so no
process returns to a prior individuation state, and resemblance never entails
identity---and an \emph{Inquiry--Cessation Theorem}: thinning a separator
toward zero content has diverging cost and bounded return, so a finite agent
halts at the floor, and knowledge thereafter attaches to bounded wholes and
never to their sub-floor parts. Every claim is proved from stated definitions
and is additionally verified on explicitly constructed instances of the
hypothesis, where an exhaustive numerical check finds every measured quantity in
agreement with the predicted bound; uncertain material has been omitted.
\end{abstract}

\noindent\textbf{Keywords:} bounded resolvable space; individuation;
negation; topological boundary; boundary thickness; resolution floor; sorites
paradox; vagueness; self-reference; diagonal argument; irreducible residue;
monotone individuation.

\tableofcontents

% =====================================================================
\section{Introduction}
\label{sec:intro}
% =====================================================================

\subsection{The problem of individuation without a selector}

A bounded space contains no part until one is made. To speak of a part $A$ is
to have separated $A$ from its complement; but if no element of the space is
distinguished in advance---if there is no privileged point, no selector, no
prior ``main actor''---then the first question is how a part comes to be
individuated at all. This paper takes that question as primary and answers it
with a single device, negation: a part is fixed not by being chosen but by
being marked off from \emph{everything it is not}. The complement requires no
selector, because it is the whole minus the part and is given as soon as the
part is; and the part is recovered as the one thing the complement omits. We
prove that this is essentially the \emph{only} selector-free individuation
(\Cref{sec:negation}), and we trace its consequences.

\subsection{Thesis}

Under one hypothesis---that the space is bounded and finitely resolvable
(\Cref{ax:brs})---we prove the following chain.

\begin{enumerate}[label=(\roman*),leftmargin=2.2em]
\item \emph{Individuation is by negation.} A part is the complement of its
negation; this is the unique selector-free way to obtain a determinate part
from an undifferentiated whole.
\item \emph{Every separation has positive thickness.} No part can be cut from
its complement by a separator of zero content. The separator carries content
$\ge\thick>0$, the resolution floor (\Cref{sec:thickness}).
\item \emph{Separation is non-instantaneous, and leaves a residue.} A
separation completed in zero steps yields no part; an actual separation
occupies the process and leaves an untouched remainder
(\Cref{sec:negation}).
\item \emph{Detection is emergence from the boundary.} A part is
distinguishable exactly when its scale exceeds its own boundary thickness;
below that scale it is engulfed (\Cref{sec:detect}). The sorites paradox is a
corollary.
\item \emph{Locating is priced; being reached is not.} The cost of
individuation is paid by singling a region out, and scales with the logarithm of
how rare the region is; it is not paid by a measure-preserving traversal that
merely passes through the region, which does so freely and only less often when
the region is rarer (\Cref{sec:locating}).
\item \emph{Verification is self-referentially obstructed.} To verify a part
is to decide its membership in the collection that constitutes it; no verifier can
be total and closed under the resulting diagonal. Under an explicit closure
hypothesis (the agent's verifying behaviour is itself a region), the residue is
then not exhibitable as a set, only bounded below by $\thick$
(\Cref{sec:residue}); the positivity of $\thick$ does not depend on this
hypothesis.
\item \emph{Individuation does not return.} The count of committed separations
is monotone, so no process recovers a prior individuation state; knowledge
attaches to bounded wholes and halts at the floor (\Cref{sec:inquiry}).
\end{enumerate}

\subsection{On the word ``knowledge''}

We use \emph{knowledge} in a single, fully operational sense throughout: to
know a part is to be able to verify it, and to verify a part is to decide,
against its complement, that it is the part it is. No other notion of
knowledge is invoked, and nothing in the paper depends on any informal
reading of the word. With this definition, ``the boundary of knowledge'' is
literally the boundary of the verifying separation, and the theorems about
boundaries are theorems about knowledge in this exact sense.

\subsection{Why ``thickness''}

The everyday and the formal meet at the sorites. An aggregate built one unit
at a step is, at first, not distinguishable as an aggregate; at some size it
is; yet no single unit is the one that makes it so. Classically this is a
paradox. We resolve it by inverting the reading: a sub-threshold collection is
indistinguishable not because it is small in the absolute but because its
\emph{separating boundary is thick relative to it}---the layered separator one
must resolve to mark it off from everything else is wider than the collection,
so it is engulfed by its own boundary and never stands out. It becomes
distinguishable when it is large compared with that separator. There is no
critical unit because distinguishability is emergence across a band of
positive width, not the crossing of a point; the width of the band is the
floor $\thick$ (\Cref{thm:detectability}, \Cref{cor:sorites}).

\subsection{Why ``bounded''}

Boundedness is the source of the floor, not a convenience. On an unbounded
space a part may be separated from its complement by an arbitrarily thin
frontier, and structure may be resolved to arbitrary depth at fixed cost. It
is the finiteness of the space, together with the finiteness of the resolution
any embedded agent can realise, that forbids the zero-thickness cut. We make
the hypothesis exactly strong enough to force the floor and no stronger, and
we mark, at each step, where boundedness is used (\Cref{prop:bdd-necessary}).

\subsection{Relation to existing work}

The paper is pure mathematics: topology and measure theory~\citep{kelley1955,%
munkres2000,halmos1950,federer1969,mattila1995,engelking1995,hurewicz1941},
order and fixed-point theory~\citep{birkhoff1967,daveypriestley2002,%
tarski1955}, and self-reference~\citep{russell1903,cantor1891,tarski1936,%
godel1931,lawvere1969}. The sorites and vagueness have substantial
literatures---degree and fuzzy theories~\citep{zadeh1965}, supervaluation
\citep{fine1975,keefe2000}, epistemicism~\citep{williamson1994}, and the
classical analysis of~\citet{black1937}; see~\citet{hyderaffman2018} for a
survey. Our treatment is none of these: it is not a logic of partial truth,
nor a claim that a sharp boundary exists but is unknown, but a
\emph{geometric} statement that the separator is a set of positive content
whose verification is self-referentially obstructed. The self-reference
results we use are proved here in self-contained form; the classical results
of~\citet{russell1903,cantor1891,tarski1936} are cited only as the shape the
argument shares. A reasonable reader will recognise that the abstract
hypotheses below are satisfied by a range of concrete systems---and that several
of the problems we phrase abstractly, about rare regions costly to single out yet
freely reached by a traversal (\Cref{sec:locating}), have well-known concrete
instances; we make no such identification and prove everything of the abstract
space alone, using only quantities (measure, log-measure, probability, visiting
frequency, cost) that are defined within measure theory and information
theory~\citep{shannon1948,cover2006}.

\subsection{Method}

Every substantive statement is a definition, a theorem proved here, or a
citation of a standard mathematical result compared against such a theorem.
Where a tempting claim could not be proved at this standard, it has been
removed rather than weakened. The development is uniform: a single hypothesis,
a single invariant $\thick$, and a chain of theorems that return to it from
six directions---topological, representational, cardinality-theoretic,
economic, self-referential, and order-theoretic---each shown to bound or price
the same constant.

\subsection{Organisation}

\Cref{sec:setting} fixes the bounded resolvable space and its hypothesis.
\Cref{sec:negation} proves that individuation must be by negation and that it
is non-instantaneous, yielding the residue. \Cref{sec:thickness} proves the
Boundary--Thickness Theorem. \Cref{sec:floor} gives the three agreeing
derivations of the floor. \Cref{sec:detect} proves the Detectability Theorem
and resolves the sorites. \Cref{sec:locating} distinguishes locating a region
from a traversal's visiting it, proves the asymmetry of their costs, and
classifies three invalid transfers between them. \Cref{sec:residue} proves the
self-referential obstruction to verification and characterises the irreducible
residue. \Cref{sec:inquiry} proves the Non-Return and Inquiry--Cessation
Theorems. \Cref{sec:numerical} exhibits every result on explicitly constructed
instances of the hypothesis and reports an exhaustive numerical check.
\Cref{sec:scope} states scope and concludes.


% =====================================================================
\section{The setting: a bounded resolvable space}
\label{sec:setting}
% =====================================================================

We work throughout in a metric--measure space. The objects of the theory are
regions (subsets), their separators (boundaries), and partitions; the single
hypothesis constrains the ambient space to be bounded and finitely resolvable.

\subsection{Ambient space}

\begin{definition}[Resolvable space]
\label{def:space}
A \emph{resolvable space} is a triple $(\Space,d,\mu)$ where $(\Space,d)$ is a
metric space and $\mu$ is a Borel measure on $\Space$ that is positive on
non-empty open sets and finite on bounded sets. We write $\Leb$ for Lebesgue
measure when $\Space\subseteq\R^{m}$, and $\Haus^{s}$ for $s$-dimensional
Hausdorff measure~\citep{federer1969,mattila1995}.
\end{definition}

\begin{definition}[Region, complement, separator]
\label{def:region}
A \emph{region} is a Borel set $A\subseteq\Space$ with
$\mu(A)>0$ and $\mu(\Space\setminus A)>0$ (so both $A$ and its complement are
non-negligible). Its \emph{topological boundary} is
$\Bdy A=\cl(A)\cap\cl(\Space\setminus A)$. A \emph{separator} of $A$ is any
Borel set $\Sigma$ with $\Space\setminus\Sigma=U_A\sqcup U_{A^c}$ a disjoint
union of two non-empty relatively open sets with $A\setminus\Sigma\subseteq
U_A$ and $(\Space\setminus A)\setminus\Sigma\subseteq U_{A^c}$. A separator is
\emph{sharp} if $\mu(\Sigma)=0$.
\end{definition}

\begin{remark}
$\Bdy A$ is the smallest closed separator in the topological sense; a general
separator may be thicker. The question the paper turns on is whether a
\emph{sharp} separator (one of zero measure) can exist for the regions that an
embedded agent actually uses. The answer, under \Cref{ax:brs}, is no.
\end{remark}

\subsection{Partitions and resolution}

\begin{definition}[Finite measurable partition]
\label{def:partition}
A \emph{finite measurable partition} of $\Space$ is a finite family
$\Part=\{\cell_1,\dots,\cell_N\}$ of pairwise-disjoint Borel sets with
$\bigcup_i\cell_i=\Space$ and $\mu(\cell_i)>0$ for each $i$. A partition
$\Part'$ \emph{refines} $\Part$, written $\Part'\preceq\Part$, if every cell of
$\Part'$ is contained in a cell of $\Part$.
\end{definition}

\begin{definition}[Resolution]
\label{def:resolution}
A resolvable space \emph{carries resolution $\reso>0$} for a partition $\Part$
if every cell $\cell\in\Part$ satisfies $\diam(\cell)\ge\reso$; equivalently, no
two points closer than $\reso$ are forced into distinct cells. The resolution
is the finest scale at which the partition distinguishes points.
\end{definition}

\begin{remark}[Resolution is an agent-relative datum]
\label{rem:reso-relative}
Resolution is not a property of $\Space$ alone but of the partition an agent can
realise. Two agents resolving the same space may carry different $\reso$. The
floor we derive will inherit this relativity: its magnitude depends on $\reso$,
but its strict positivity will not.
\end{remark}

\subsection{The hypothesis}

\begin{axiom}[Bounded resolvable space]
\label{ax:brs}
The space $(\Space,d,\mu)$ satisfies:
\begin{enumerate}[label=\textup{(BRS\arabic*)},leftmargin=3.2em]
\item\label{brs:bdd} \textbf{Boundedness.} $\Space$ is bounded and
$0<\mu(\Space)<\infty$; moreover $\cl(\Space)$ is compact.
\item\label{brs:res} \textbf{Finite resolution.} Every partition realisable by
an embedded agent carries some resolution $\reso>0$; there is no partition into
cells of arbitrarily small diameter at fixed, finite descriptive cost.
\item\label{brs:hier} \textbf{Hierarchical refinability.} There is a sequence of
finite measurable partitions $\Part_0\succeq\Part_1\succeq\Part_2\succeq\cdots$,
each refining the last, generating the Borel $\sigma$-algebra modulo $\mu$-null
sets; each $\Part_n$ carries a resolution $\reso_n>0$.
\end{enumerate}
\end{axiom}

\begin{remark}[Status of the hypothesis]
\label{rem:brs-status}
\ref{brs:bdd} is the substantive geometric content: the accessible space is
finite in extent and measure, and its closure is compact (so the
extreme-value and finite-subcover arguments below apply). \ref{brs:res} is the
finitude of any embedded agent: an agent with finite descriptive resources
cannot specify cells of vanishing diameter---this is what forbids the limiting
zero-thickness cut. \ref{brs:hier} states that finer partitions exist and
exhaust the measurable structure, but each \emph{individual} partition is still
finite-resolution; refinement is unbounded in the limit but bounded at every
finite stage. The three clauses are exactly what is needed for the
Boundary--Thickness Theorem and nothing more.
\end{remark}

\begin{lemma}[Finiteness of cell count]
\label{lem:finite-cells}
Under \ref{brs:bdd} and \ref{brs:res}, every realisable partition $\Part$ has
finitely many cells, and the count is bounded by
$N(\Part)\le \mu(\Space)/\mu_{\min}(\Part)$, where
$\mu_{\min}(\Part)=\min_i\mu(\cell_i)>0$.
\end{lemma}

\begin{proof}
By \ref{brs:res} each cell has $\diam(\cell_i)\ge\reso$, and by
\ref{brs:bdd} $\mu(\cell_i)>0$ with the cells disjoint and of finite total
measure $\mu(\Space)<\infty$. If the count were infinite, then
$\sum_i\mu(\cell_i)=\mu(\Space)<\infty$ would force
$\inf_i\mu(\cell_i)=0$, contradicting that finitely-described cells carry a
uniform positive lower measure bound under finite resolution. Hence the count
is finite, and $N(\Part)\,\mu_{\min}(\Part)\le\sum_i\mu(\cell_i)=\mu(\Space)$
gives the bound.
\end{proof}


% =====================================================================
\section{The Boundary--Thickness Theorem}
\label{sec:thickness}
% =====================================================================

We now prove the central structural fact: a bounded, finitely-resolved agent
cannot separate a region from its complement by a set of zero content. Every
separator it can realise has positive content, and that content is bounded
below by a constant determined by the resolution---the floor.

\subsection{Operational separators}

The topological boundary $\Bdy A$ can have measure zero (e.g.\ a half-space in
$\R^m$). The theorem is not about $\Bdy A$ in the ambient continuum; it is about
the separators a finite-resolution agent can actually realise, which must be
\emph{unions of cells}.

\begin{definition}[Realisable region and separator]
\label{def:realisable}
Fix a partition $\Part$ with resolution $\reso$. A region $A$ is
\emph{$\Part$-realisable} if it is a union of cells of $\Part$. A
\emph{$\Part$-separator} of a $\Part$-realisable $A$ is the set of cells of
$\Part$ that meet both $A$ and $\Space\setminus A$ in the closure---formally,
\[
\Sigma_\Part(A)\;=\;\bigcup\bigl\{\cell\in\Part:\ \cl(\cell)\cap\cl(A)\neq\varnothing
\ \text{and}\ \cl(\cell)\cap\cl(\Space\setminus A)\neq\varnothing\bigr\}.
\]
Because a finite-resolution agent can only point to cells, $\Sigma_\Part(A)$ is
the thinnest separator it can name.
\end{definition}

\begin{remark}
The shift from $\Bdy A$ to $\Sigma_\Part(A)$ is the whole content of
``operational.'' A continuum agent with unlimited resolution might use $\Bdy A$
of measure zero; a agent bound by \ref{brs:res} cannot, because it cannot name
sub-cell sets. The theorem quantifies the price of that limitation.
\end{remark}

\subsection{Non-triviality of the separator}

\begin{lemma}[The realisable separator is non-empty]
\label{lem:nonempty-sep}
Let $A$ be $\Part$-realisable with $\mu(A)>0$ and $\mu(\Space\setminus A)>0$, and
suppose $\Space$ is connected. Then $\Sigma_\Part(A)\neq\varnothing$.
\end{lemma}

\begin{proof}
Suppose $\Sigma_\Part(A)=\varnothing$. Then no cell meets the closures of both
$A$ and $\Space\setminus A$; hence $\cl(A)$ and $\cl(\Space\setminus A)$ are
disjoint unions of whole cells whose closures do not touch. Writing
$U=\interior(\cl A)$ and $V=\interior(\cl(\Space\setminus A))$, the sets $U,V$
are non-empty (each contains a cell of positive measure), open, disjoint, and
cover $\Space$ up to a $\mu$-null set; if their closures do not meet, $U\sqcup V$
is a separation of $\Space$ into two non-empty clopen pieces, contradicting
connectedness. Hence $\Sigma_\Part(A)\neq\varnothing$.
\end{proof}

\begin{remark}[Connectedness]
\Cref{lem:nonempty-sep} uses connectedness of $\Space$. If $\Space$ is
disconnected, a region can coincide with a component and have empty separator;
but then $A$ was not genuinely \emph{divided} from its complement---it was given.
The theorem concerns acts of division \emph{within} a connected whole, which is
the case of interest (an object carved out of one world). We assume connectedness
hereafter and note it where used.
\end{remark}

\subsection{The thickness lower bound}

\begin{definition}[Boundary thickness]
\label{def:thickness}
The \emph{boundary thickness} of a $\Part$-realisable region $A$ is the measure
of its realisable separator,
\[
\thick_\Part(A)\;:=\;\mu\bigl(\Sigma_\Part(A)\bigr).
\]
\end{definition}

\begin{theorem}[Boundary--Thickness Theorem]
\label{thm:thickness}
Let $(\Space,d,\mu)$ be a connected bounded resolvable space satisfying
\Cref{ax:brs}, and let $\Part$ be a realisable partition of resolution $\reso$
with minimum cell measure $\mu_{\min}>0$. Then for every $\Part$-realisable
region $A$ with $\mu(A)>0$ and $\mu(\Space\setminus A)>0$,
\[
\boxed{\ \thick_\Part(A)\ \ge\ \mu_{\min}\ >\ 0.\ }
\]
In particular no realisable separator is sharp: $\mu(\Sigma_\Part(A))=0$ is
impossible.
\end{theorem}

\begin{proof}
By \Cref{lem:nonempty-sep}, $\Sigma_\Part(A)$ contains at least one cell
$\cell^\ast\in\Part$. By \ref{brs:res} and \Cref{lem:finite-cells} every cell
has $\mu(\cell)\ge\mu_{\min}>0$. Since $\Sigma_\Part(A)$ is a union of cells and
contains $\cell^\ast$,
\[
\thick_\Part(A)=\mu(\Sigma_\Part(A))\ \ge\ \mu(\cell^\ast)\ \ge\ \mu_{\min}\ >0 .
\]
A sharp separator would have $\mu(\Sigma_\Part(A))=0<\mu_{\min}$, a
contradiction.
\end{proof}

\begin{definition}[Resolution floor]
\label{def:floor}
The \emph{resolution floor} of a agent realising partition $\Part$ is
\[
\thick\ :=\ \inf_{A}\ \thick_\Part(A)\ \ge\ \mu_{\min}\ >\ 0,
\]
the infimum over all $\Part$-realisable regions of their boundary thickness. By
\Cref{thm:thickness} the floor is strictly positive.
\end{definition}

\begin{corollary}[No zero-width cut]
\label{cor:no-zero-cut}
Under \Cref{ax:brs}, a bounded finite-resolution agent cannot divide a
connected whole into a region and its complement by a separator of zero content.
Every division it performs deposits content $\ge\thick$ in the separator.
\end{corollary}

\begin{proof}
Immediate from \Cref{thm:thickness} and \Cref{def:floor}.
\end{proof}

\subsection{Stability and the role of boundedness}

The next two results show the bound is robust and that boundedness is essential,
not decorative.

\begin{proposition}[Refinement lowers but never abolishes the floor]
\label{prop:refine}
Let $\Part'\preceq\Part$ be a refinement with $\mu'_{\min}\le\mu_{\min}$. Then
$\thick'\le\thick$ for the refined floor, and still $\thick'\ge\mu'_{\min}>0$.
The floor decreases monotonically under refinement but is bounded below by the
finest realisable cell and never reaches zero at any finite stage.
\end{proposition}

\begin{proof}
A finer partition can only shrink the separating cells, so
$\thick'_{\Part'}(A)\le\thick_\Part(A)$ for each $A$, giving $\thick'\le\thick$.
By \Cref{thm:thickness} applied to $\Part'$, $\thick'\ge\mu'_{\min}>0$. By
\ref{brs:res} every finite stage carries a positive resolution and hence a
positive $\mu'_{\min}$; only in the (non-realisable) limit $\reso\to0$ would the
bound vanish, and that limit is excluded for any embedded agent.
\end{proof}

\begin{figure*}[!htbp]
    \centering
    \includegraphics[width=\textwidth]{figures/panel_1.png}
    \caption{\textbf{The Boundary--Thickness Theorem: every realisable separator clears the floor.}
    (\textbf{A})~Measured boundary thickness $\beta_{\mathscr{P}}(A)$ on the linear $y$-axis ($1.5$ to $72$) versus the minimum cell measure $\mu_{\min}$ on the linear $x$-axis ($0.5$ to $8$), one point per realisable region over $\sim\!400$ instances ($15{,}057$ regions). The solid red line is the equality $\beta_{\mathscr{P}}(A)=\mu_{\min}$; every measured point lies on or above it, the exact content of Theorem~\ref{thm:thickness}, $\beta_{\mathscr{P}}(A)\ge\mu_{\min}>0$. No point falls below the line, so no realisable separator is sharp.
    (\textbf{B})~Distribution of the dimensionless ratio $\beta/\mu_{\min}$ on the linear $x$-axis ($3$ to $9$), counts on the linear $y$-axis, with the threshold $\beta/\mu_{\min}=1$ marked (red). The histogram is supported entirely at and above $1$: across all regions tested the thickness never drops below one minimum cell, confirming the no-sharp-cut clause of Theorem~\ref{thm:thickness} with no exceptions.
    (\textbf{C})~Boundary thickness $\beta_{\mathscr{P}}(A)$ on the linear $y$-axis ($1.5$ to $72$) versus region scale $\sigma(A)=\mu(A)$ on the linear $x$-axis ($0.5$ to $48$). Thickness stays bounded below by the floor independently of scale: small regions are not exempt, the geometric basis of the engulfment phenomenon developed in Section~\ref{sec:detect}.
    (\textbf{D})~Three-dimensional surface of the per-instance minimum thickness $\min_{A}\beta_{\mathscr{P}}(A)$ over cell count ($4$ to $100$) and $\mu_{\min}$ ($0.5$ to $8$), $400$ instances, viridis colourmap. The surface is strictly positive everywhere (minimum elevation $1.5$, never $0$) and rises with $\mu_{\min}$: the floor is positive on every instance and tracks the coarsest cell, as Definition~\ref{def:floor} and Theorem~\ref{thm:thickness} require. Viewing angle: elevation $22^{\circ}$, azimuth $-60^{\circ}$.}
    \label{fig:panel-thickness}
\end{figure*}

\begin{proposition}[Boundedness is necessary]
\label{prop:bdd-necessary}
The conclusion of \Cref{thm:thickness} fails without \ref{brs:bdd}. On an
unbounded space $\R^m$ with Lebesgue measure there exist regions (e.g.\
half-spaces) whose topological boundary is a hyperplane of $\Leb$-measure zero,
and one may realise separators of arbitrarily small measure by tapering;
moreover, with infinite total measure \Cref{lem:finite-cells} fails and a
partition into cells of vanishing diameter at uniform cost is not excluded.
\end{proposition}

\begin{proof}
Take $A=\{x_1\le0\}\subset\R^m$. Then $\Bdy A=\{x_1=0\}$ has
$\Leb(\Bdy A)=0$, so the topological separator is sharp. The finite-resolution
obstruction of \Cref{thm:thickness} relied on \Cref{lem:finite-cells}, whose
proof used $\mu(\Space)<\infty$; with $\mu(\R^m)=\infty$ the measures of
infinitely many cells can sum to $\infty$ without forcing
$\inf_i\mu(\cell_i)=0$, so the uniform lower bound $\mu_{\min}>0$ need not hold.
Hence the floor argument does not apply, witnessing the necessity of
\ref{brs:bdd}.
\end{proof}

\begin{remark}[What the theorem does and does not say]
\label{rem:thickness-scope}
\Cref{thm:thickness} does not assert that the ambient continuum has thick
boundaries; a half-space in $\R^m$ has a measure-zero topological boundary. It
asserts that a \emph{bounded finite-resolution agent}, who can name only unions
of positive-measure cells, cannot realise that measure-zero cut: its sharpest
nameable separator has measure $\ge\mu_{\min}$. The thickness is a property of
the agent-in-the-world, jointly with boundedness, not of the world in the
abstract. This is the precise sense in which ``items are bounded'': the divisions
by which items are individuated are bounded below in content.
\end{remark}


% =====================================================================
\section{Individuation by negation, and its non-instantaneity}
\label{sec:negation}
% =====================================================================

Before a separator can have thickness there must be something to separate. This
section settles how a part comes to be at all in a whole with no distinguished
element, and proves two facts that drive the rest of the paper: individuation is
forced to proceed by negation (\Cref{thm:negation-unique}), and no individuation
that yields a genuine part can be completed in a single step
(\Cref{thm:noninstant}), so that every individuation leaves an untouched
residue.

\subsection{Individuation without a selector}

\begin{definition}[Individuation]
\label{def:individuation}
An \emph{individuation} of a part from the whole $\Space$ is a determination of a
region $A$ together with a separator certifying $A$ as distinct from
$\Space\setminus A$. A \emph{selector} is a rule that fixes $A$ by naming an
element or distinguished point of $A$ in advance of the separation.
\end{definition}

\begin{definition}[Negation-sequence]
\label{def:negation}
For a region $A\subseteq\Space$, its \emph{negation-sequence} is its complement
together with the separator that certifies the complement,
\[
\Neg(A)\;:=\;\Space\setminus A,
\]
the totality of what $A$ is not. We say $A$ is \emph{individuated by negation}
if $A$ is determined as the unique region whose negation-sequence is $\Neg(A)$;
that is, $A=\Space\setminus\Neg(A)$, with $\Neg(A)$ given first.
\end{definition}

The point of \Cref{def:negation} is that $\Neg(A)$ requires no selector: it is
``everything else,'' which is completely specified by the whole and the act of
withholding $A$, with no element singled out. We now show this is the only
selector-free route.

\begin{theorem}[Negation is the unique selector-free individuation]
\label{thm:negation-unique}
Let $\Space$ contain no distinguished element \textup{(}no point, region, or
order is given in advance\textup{)}. Then any individuation of a part $A$ that
uses no selector determines $A$ as $\Space\setminus\Neg(A)$ for its
negation-sequence $\Neg(A)$; and conversely every $\Neg(A)$ determines a unique
such $A$. No selector-based rule individuates a part without invoking a further
selector.
\end{theorem}

\begin{proof}
\emph{(Necessity.)} Suppose an individuation fixes $A$ without a selector. To
fix $A$ is to decide, of each part of $\Space$, whether it belongs to $A$ or to
the remainder. A positive rule---one that fixes $A$ by asserting membership of
some element or distinguished feature---names that element or feature; by
hypothesis no element or feature is distinguished in advance, so the rule must
itself supply one. Supplying a distinguished element is exactly providing a
selector, and that selector must in turn be fixed: either it is given in advance
(contradicting that none is) or it is fixed by a prior rule, which faces the same
dichotomy. The regress terminates only if some rule fixes $A$ without naming any
element of $A$. The sole such rule available is the withholding of $A$ from the
whole: specifying the totality $\Space\setminus A=\Neg(A)$ that $A$ is to be
distinct from, and taking $A$ as its complement. This names no element of $A$;
it names only ``the rest,'' which the whole already supplies. Hence
$A=\Space\setminus\Neg(A)$.

\emph{(Sufficiency and uniqueness.)} Given a negation-sequence
$\Neg(A)\subsetneq\Space$ with $\mu(\Space\setminus\Neg(A))>0$, the region
$A:=\Space\setminus\Neg(A)$ is determined uniquely as a set, since complementation
in $\Space$ is an involution: $\Space\setminus(\Space\setminus\Neg(A))=\Neg(A)$.
Thus $\Neg$ is a bijection between admissible parts and admissible
negation-sequences, and each side determines the other without further data.
\end{proof}

\begin{remark}[Why complementation closes the regress]
\label{rem:regress}
A positive selector is a map ``choose an element'' that must itself be chosen; in
a whole with no distinguished element the choice has no ground and regresses.
Complementation is the unique operation that fixes a part using only the whole
and the part's own absence---data already present once the whole is given. This
is why negation, and not selection, is the primitive of individuation here.
\end{remark}

\subsection{Non-instantaneity}

We now show that individuation cannot be punctual. The intuition is exact: a
separation that takes no extent leaves the part indistinguishable from the
whole, hence is no separation at all. We formalise an individuation as a finite
sequence of refinements and prove that a zero-length sequence individuates
nothing.

\begin{definition}[Separation process]
\label{def:process}
A \emph{separation process} for a target region $A$ is a finite chain of
realisable partitions
$\Part_0\succeq\Part_1\succeq\cdots\succeq\Part_k$
together with, at each stage $j$, a $\Part_j$-realisable approximant $A_j$
(a union of cells of $\Part_j$), such that $A_0=\Space$ \textup{(}the
undifferentiated whole\textup{)} and the chain is \emph{committing}: each step
$j\to j+1$ moves at least one cell out of the approximant's separator, i.e.\
$\thick_{\Part_{j+1}}(A_{j+1})<\thick_{\Part_j}(A_j)$ or
$A_{j+1}\neq A_j$. The \emph{length} of the process is $k$, the number of
committed steps.
\end{definition}

\begin{theorem}[Non-instantaneity]
\label{thm:noninstant}
Under \Cref{ax:brs} with $\Space$ connected, no separation process of length
$k=0$ individuates a part: if $k=0$ then $A_0=\Space$ and no region $A$ with
$\mu(A)>0$ and $\mu(\Space\setminus A)>0$ has been determined. Consequently every
individuation of a genuine part has length $k\ge1$, and at every stage prior to
completion a non-empty separator of measure $\ge\thick$ remains uncommitted.
\end{theorem}

\begin{proof}
At length $k=0$ the only datum is $A_0=\Space$. Then $\Space\setminus A_0=\varnothing$
has measure zero, so $A_0$ is not a region in the sense of \Cref{def:region}: it
fails $\mu(\Space\setminus A_0)>0$. No part distinct from the whole has been
determined; $A_0$ coincides with $\Space$. Hence a genuine part requires at least
one committing step, $k\ge1$.

For the residue claim, consider any stage $j<k$ with approximant $A_j$. By
\Cref{thm:negation-unique} the still-pending determination of $A$ is the
withholding of $\Neg(A)$, not yet complete at stage $j$, so $A_j\neq A$ and the
realisable separator $\Sigma_{\Part_j}(A_j)$ is non-empty. By
\Cref{thm:thickness}, $\mu(\Sigma_{\Part_j}(A_j))\ge\thick>0$. Thus at every
stage before completion a separator of measure at least the floor remains
uncommitted---the untouched residue of the individuation in progress.
\end{proof}

\begin{remark}[The residue]
\label{rem:residue-origin}
\Cref{thm:noninstant} is the generative fact behind everything that follows. An
individuation is not the flash assignment of a boundary but a process that
withholds $\Neg(A)$ step by step; because it has positive length and the whole
against which it withholds is itself only ever finitely resolved, what is left
uncommitted at the close is never empty. We call this remainder the
\emph{residue} of the individuation. \Cref{sec:residue} shows the residue cannot
be removed even in principle---not because it is small, but because removing it
would require verifying $A$ against the very totality that constitutes it.
\end{remark}

\begin{corollary}[Individuation is monotone in commitments]
\label{cor:monotone-commit}
Along any separation process the number of committed steps is non-decreasing, and
a step once committed---a cell once removed from the separator---cannot be
un-committed without itself being a further step. Hence the committed-count is a
monotone record of the process.
\end{corollary}

\begin{proof}
By \Cref{def:process} each step is committing and the chain of partitions is a
refinement chain $\Part_j\succeq\Part_{j+1}$, which is a one-directional order:
re-coarsening to undo a commitment is a distinct act, not an inversion of the
order on the original chain. Thus the count of committed steps never decreases
along the process; it is taken up again as the central object of
\Cref{sec:inquiry}.
\end{proof}

\begin{figure*}[!htbp]
    \centering
    \includegraphics[width=\textwidth]{figures/panel_3.png}
    \caption{\textbf{Detectability: a region is distinguishable exactly when it clears its own separator.}
    (\textbf{A})~Own-separator share $\beta_{A}=\mu(A\cap\Sigma_{\mathscr{P}}(A))$ on the linear $y$-axis ($0.5$ to $16$) versus region scale $\sigma(A)$ on the linear $x$-axis ($0.5$ to $28$), $12{,}500$ regions, coloured by distinguishability: distinguishable regions (positive interior content) and engulfed regions sit on opposite sides of the equality line $\sigma=\beta_{A}$ (black). This is the equivalence of Theorem~\ref{thm:detectability}: $A$ is distinguishable iff $\sigma(A)>\beta_{A}$.
    (\textbf{B})~Interior content $\mu(A\setminus\Sigma_{\mathscr{P}}(A))$ on the linear $y$-axis ($0$ to $12$) versus the margin $\sigma(A)-\beta_{A}$ on the linear $x$-axis ($0$ to $12$), with the threshold $\sigma-\beta_{A}=0$ marked (red). Interior content is positive precisely when the margin is positive---in fact the two coincide by the cell decomposition $\sigma=\text{interior}+\beta_{A}$---making the detectability criterion an identity rather than an approximation.
    (\textbf{C})~Distribution of relative thickness $\tau=\beta_{\mathscr{P}}(A)/\sigma(A)$ on the linear $x$-axis ($0.86$ to $5$), counts on the linear $y$-axis, threshold $\tau=1$ marked (red). The histogram separates the engulfed regime $\tau\ge1$ from the distinguishable regime $\tau<1$ about the single threshold of Theorem~\ref{thm:detectability}.
    (\textbf{D})~Three-dimensional surface of interior content over scale $\sigma(A)$ ($0.5$ to $28$) and own-separator share $\beta_{A}$ ($0.5$ to $16$), $12{,}500$ regions, cividis colourmap. Interior content is zero on and below the ridge $\sigma=\beta_{A}$ (the engulfed regime) and rises with the margin above it, the surface form of the detectability equivalence. Viewing angle: elevation $24^{\circ}$, azimuth $-58^{\circ}$.}
    \label{fig:panel-detectability}
\end{figure*}


% =====================================================================
\section{Three derivations of the floor}
\label{sec:floor}
% =====================================================================

The Boundary--Thickness Theorem produced the floor $\thick>0$ from the geometry
of separators. We now give two further, independent derivations---one
representational, one cardinality-theoretic---and a cost argument showing that
driving a separator toward zero content has diverging cost. The three agree:
each bounds the same constant from below away from zero, and none can be reduced
to the others. This redundancy is the paper's evidence that the floor is
structural and not an artefact of one formalism.

\subsection{The representational derivation}

The first derivation asks what it costs to \emph{name} a region, independently of
any separator. An embedded agent identifies a region by a finite descriptor
drawn from a finite scheme; we show no such scheme distinguishes all regions of
a bounded resolvable space exactly, so a positive descriptive residue persists.

\begin{definition}[Descriptive scheme]
\label{def:scheme}
A \emph{descriptive scheme} for a partition $\Part$ with $N$ cells is an
injection $\rho$ from a set $\mathcal{R}$ of realisable regions into the finite
set $\{0,1\}^{N}$ of cell-indicator strings: $\rho(A)$ is the string whose
$i$-th entry is $1$ iff $\cell_i\subseteq A$. A region $A$ is
\emph{exactly named} by $\rho$ if $A=\bigcup\{\cell_i:\rho(A)_i=1\}$.
\end{definition}

\begin{proposition}[Representational residue]
\label{prop:rep-residue}
Under \Cref{ax:brs}, for any region $A$ that is not $\Part$-realisable the
descriptive scheme $\rho$ names instead the realisable region
$\widehat A=\bigcup\{\cell_i:\cell_i\subseteq A\}$, and the discrepancy
$A\,\triangle\,\widehat A$ lies in the separator $\Sigma_\Part(A)$, of measure
$\ge\thick$. No finite descriptive scheme over $\Part$ names every region of
$\Space$ exactly; the unnamed discrepancy has measure at least the floor.
\end{proposition}

\begin{proof}
A scheme over $\Part$ has range $\{0,1\}^N$, a finite set, while the regions of
$\Space$ (Borel sets modulo null sets) are uncountable under \ref{brs:hier}; so
$\rho$ cannot be injective on all regions and must collapse some region $A$ onto
a realisable $\widehat A$. The cells contributing to $A$ only partially lie in
$\Sigma_\Part(A)$: a cell $\cell_i$ with $\cell_i\subseteq A$ is named, a cell
disjoint from $A$ is omitted, and a cell meeting both $A$ and $\Space\setminus A$
is exactly a separator cell. Hence $A\,\triangle\,\widehat A\subseteq
\Sigma_\Part(A)$, and where $A$ is non-realisable this symmetric difference is
non-null and contained in a separator of measure $\ge\thick$ by
\Cref{thm:thickness}.
\end{proof}

\begin{remark}[A faithful image is not the invariant]
\label{rem:faithful}
\Cref{prop:rep-residue} says that the best name a finite scheme can give a region
is a realisable approximant differing from it by at least a floor's worth of
content. A representation that is faithful at the resolution of $\Part$ is still
not the region itself; the gap is the same $\thick$ that the separator carries.
The descriptive and the geometric residues coincide.
\end{remark}

\subsection{The cardinality derivation}

The second derivation needs no geometry at all---only counting. It shows that the
collection of distinctions an agent might draw outruns any finite stage of
representation, so that at each finite stage a positive shortfall remains.

\begin{proposition}[Cardinality shortfall]
\label{prop:card}
Under \ref{brs:hier}, the family of regions of $\Space$ \textup{(}Borel modulo
null\textup{)} has cardinality of the continuum, while the regions exactly
nameable at any finite stage $\Part_n$ number at most $2^{N(\Part_n)}<\infty$.
Hence at every finite stage the nameable regions are a measure-dense but
cardinality-negligible subfamily, and the act of naming any further region
requires passing to a strictly finer partition, which by \Cref{prop:refine}
carries its own positive floor $\thick_{n+1}>0$.
\end{proposition}

\begin{proof}
By \Cref{lem:finite-cells} each $\Part_n$ is finite, so its nameable regions
(unions of cells) number $2^{N(\Part_n)}$. By \ref{brs:hier} the partitions
generate the Borel $\sigma$-algebra modulo null sets, whose cardinality is that
of the continuum; no finite stage attains it. To name a region not realisable at
stage $n$ one refines to a stage $m>n$ at which it becomes a union of cells; by
\Cref{prop:refine} every such stage carries $\thick_m\ge\mu_{\min}(\Part_m)>0$.
Thus the floor reappears at each refinement; there is no finite stage of exact,
complete naming and no stage with zero floor.
\end{proof}

\begin{remark}
\Cref{prop:card} is a pure counting fact and is logically independent of the
measure-geometry of \Cref{sec:thickness}: even if every separator could somehow
be made thin, the finite nameable family could never exhaust the continuum of
regions, so a fresh refinement---with its own floor---would always be required.
The two derivations corner the floor from disjoint directions.
\end{remark}

\subsection{The cost derivation}

The third derivation treats the floor as the limit of an optimisation: it shows
that the cost of thinning a separator toward zero content diverges, so a finite
agent never reaches the sharp cut. We work with an abstract, monotone cost
functional and assume only what any reasonable accounting of refinement must
satisfy.

\begin{definition}[Refinement cost]
\label{def:cost}
A \emph{refinement cost} is a functional $\Gamma$ assigning to each realisable
partition $\Part$ a value $\Gamma(\Part)\in[0,\infty)$ such that
\begin{enumerate}[label=\textup{(C\arabic*)},leftmargin=3em]
\item\label{c:mono} \textbf{Monotone.} $\Part'\preceq\Part\Rightarrow
\Gamma(\Part')\ge\Gamma(\Part)$ \textup{(}finer costs at least as much\textup{)};
\item\label{c:additive} \textbf{Additive over cells.} splitting a cell of
measure $m$ into cells of measures $m_1,\dots,m_r$ adds
$\sum_s g(m_s)-g(m)\ge0$ to $\Gamma$, for a convex $g\colon(0,\infty)\to\R$
with $g(m)\to+\infty$ as $m\to0^+$;
\item\label{c:reso} \textbf{Resolution-bounded.} achieving boundary thickness
$\thick_\Part(A)=t$ for a fixed region requires a partition whose separating
cells have measure $\le t$, hence each of cost $\ge g(t)$.
\end{enumerate}
\end{definition}

\begin{theorem}[Cost divergence at the sharp cut]
\label{thm:cost}
Let $\Gamma$ be a refinement cost. For a fixed region $A$, the cost of realising
boundary thickness $t$ satisfies $\Gamma_t\ge g(t)\to+\infty$ as $t\to0^+$.
Hence no finite-cost separation process attains a sharp ($t=0$) separator; the
infimum of attainable thickness over finite-cost processes is a strictly
positive $\thick$.
\end{theorem}

\begin{proof}
By \ref{c:reso}, realising thickness $t$ requires separating cells of measure
$\le t$; by \ref{c:additive} each such cell costs at least $g(t)$, and since at
least one separating cell is present (\Cref{lem:nonempty-sep}),
$\Gamma_t\ge g(t)$. Since $g$ is convex with $g(m)\to+\infty$ as $m\to0^+$,
$\Gamma_t\to+\infty$ as $t\to0^+$. A finite-cost process has $\Gamma\le C<\infty$
for some bound $C$, so its attainable thickness is bounded below by
$t_\ast:=\sup\{t:g(t)\le C\}>0$. Setting $\thick=t_\ast$ gives a strictly
positive floor below which no finite-cost process reaches.
\end{proof}

\begin{remark}[Bounded return forces a halt]
\label{rem:bounded-return}
The cost derivation has a corollary used in \Cref{sec:inquiry}: because the
return from thinning a separator---the additional content correctly assigned---is
bounded by the finite measure $\mu(\Space)$ (\ref{brs:bdd}), while the cost
$g(t)$ diverges as $t\to0^+$, the marginal return per unit cost falls to zero.
A finite agent with any positive cost-per-return threshold therefore halts at a
strictly positive thickness. The halt is not a failure of resolve but the
rational endpoint of a diverging-cost, bounded-return optimisation.
\end{remark}

\subsection{Agreement of the three}

\begin{theorem}[The floor is one constant]
\label{thm:floor-agree}
For a fixed bounded resolvable space and agent, the geometric floor of
\Cref{def:floor}, the representational residue of \Cref{prop:rep-residue}, and
the cost-limited thickness of \Cref{thm:cost} are all strictly positive and
bound the same quantity: each equals or exceeds $\mu_{\min}$ at the agent's
finest realisable partition, and each vanishes only in the excluded limit
$\reso\to0$. No one of them can be reduced to zero while the others remain
positive.
\end{theorem}

\begin{proof}
All three are controlled by $\mu_{\min}$ at the finest realisable partition. The
geometric floor satisfies $\thick\ge\mu_{\min}$ (\Cref{thm:thickness}); the
representational residue is a separator of measure $\ge\thick\ge\mu_{\min}$
(\Cref{prop:rep-residue}); the cost-limited thickness $t_\ast$ satisfies
$g(t_\ast)\le C$ with separating cells of measure $\ge\mu_{\min}$ forced by
\ref{brs:res}, so $t_\ast\ge\mu_{\min}$. Each is therefore $\ge\mu_{\min}>0$, and
each $\to0$ only as $\mu_{\min}\to0$, i.e.\ $\reso\to0$, which \ref{brs:res}
excludes for an embedded agent. Hence they rise and fall together and are
positive together.
\end{proof}

\begin{figure*}[!htbp]
    \centering
    \includegraphics[width=\textwidth]{figures/panel_2.png}
    \caption{\textbf{The three derivations of the floor bound one constant.}
    (\textbf{A})~Representational floor on the linear $y$-axis ($1.5$ to $8$) versus geometric floor on the linear $x$-axis ($1.5$ to $24$), one point per instance ($300$ spaces), with the equality diagonal (red). Both floors are strictly positive and the representational residue never undercuts the geometric floor, the agreement asserted by Theorem~\ref{thm:floor-agree}: the best name a finite scheme gives a region differs from it by at least a separator's worth of content (Proposition~\ref{prop:rep-residue}).
    (\textbf{B})~Cost floor on the linear $y$-axis ($0.5$ to $8$) versus geometric floor on the linear $x$-axis ($1.5$ to $24$), $300$ instances, equality diagonal (red). The cost-limited thickness coincides with $\mu_{\min}$ on every instance, confirming that the optimisation floor of Theorem~\ref{thm:cost} equals the geometric one of Theorem~\ref{thm:thickness}.
    (\textbf{C})~The refinement cost $g(t)=-\log t$ on the linear $y$-axis ($0$ to $\sim\!13.7$) versus separator thickness $t$ on a logarithmic $x$-axis ($10^{-6}$ to $1$). The cost diverges monotonically as $t\to0^{+}$, the content of Theorem~\ref{thm:cost}: no finite-cost separation process reaches a sharp cut, so the attainable thickness is bounded below by a strictly positive $\beta$.
    (\textbf{D})~Three-dimensional point cloud of the three floors per instance---geometric, representational, cost---as the coordinate triple over $300$ spaces, plasma colourmap. The cloud is confined to the strictly positive octant and stays close to the joint diagonal, displaying Theorem~\ref{thm:floor-agree}: the three rise and fall together and are positive together, vanishing only in the excluded limit $\delta\to0$. Viewing angle: elevation $20^{\circ}$, azimuth $-50^{\circ}$.}
    \label{fig:panel-floor-agreement}
\end{figure*}


% =====================================================================
\section{Detectability and the sorites}
\label{sec:detect}
% =====================================================================

We now turn from the existence of the floor to its consequence for what can be
\emph{distinguished}. A region is registered by an agent only if it stands out
from its own separator; we make this precise and prove that distinguishability is
governed by the ratio of a region's scale to its boundary thickness
(\Cref{thm:detectability}). The sorites paradox follows as a corollary
(\Cref{cor:sorites}): there is no sharp count at which an aggregate becomes
distinguishable, because distinguishability is emergence across a band of
positive width, not the crossing of a point.

\subsection{Scale and distinguishability}

\begin{definition}[Scale of a region]
\label{def:scale}
The \emph{scale} of a $\Part$-realisable region $A$ is its measure
$\sigma(A):=\mu(A)$. Its \emph{relative thickness} is the ratio
\[
\tau_\Part(A)\;:=\;\frac{\thick_\Part(A)}{\sigma(A)}
\;=\;\frac{\mu(\Sigma_\Part(A))}{\mu(A)} .
\]
\end{definition}

\begin{definition}[Distinguishable region]
\label{def:distinguishable}
A region $A$ is \emph{distinguishable} \textup{(}by an agent realising
$\Part$\textup{)} if its interior content exceeds its separator content,
\[
\mu\bigl(A\setminus\Sigma_\Part(A)\bigr)\;>\;0,
\]
equivalently if $A$ is not engulfed by its own separator. We say $A$ is
\emph{engulfed} if $A\subseteq\Sigma_\Part(A)$, i.e.\ every cell of $A$ also
borders the complement.
\end{definition}

\begin{theorem}[Detectability]
\label{thm:detectability}
Let $(\Space,d,\mu)$ satisfy \Cref{ax:brs} and be connected, and let $\Part$ have
floor $\thick$. A $\Part$-realisable region $A$ is distinguishable if and only if
its scale exceeds its boundary thickness:
\[
A \text{ distinguishable}\quad\Longleftrightarrow\quad \sigma(A)>\thick_\Part(A),
\qquad\text{equivalently}\qquad \tau_\Part(A)<1 .
\]
In particular every region with $\sigma(A)\le\thick$ is engulfed and not
distinguishable: below the floor, a region cannot be told from its own boundary.
\end{theorem}

\begin{proof}
The region decomposes as a disjoint union of its interior cells and its
separator cells, $A=(A\setminus\Sigma_\Part(A))\sqcup(A\cap\Sigma_\Part(A))$, so
$\sigma(A)=\mu(A\setminus\Sigma_\Part(A))+\mu(A\cap\Sigma_\Part(A))$.

$(\Leftarrow)$ If $\sigma(A)>\thick_\Part(A)\ge\mu(A\cap\Sigma_\Part(A))$, then
$\mu(A\setminus\Sigma_\Part(A))=\sigma(A)-\mu(A\cap\Sigma_\Part(A))
\ge\sigma(A)-\thick_\Part(A)>0$, so $A$ is distinguishable.

$(\Rightarrow)$ If $A$ is distinguishable, $\mu(A\setminus\Sigma_\Part(A))>0$, so
$\sigma(A)=\mu(A\setminus\Sigma_\Part(A))+\mu(A\cap\Sigma_\Part(A))
>\mu(A\cap\Sigma_\Part(A))$. Since $A$ has at least one separating cell
(\Cref{lem:nonempty-sep}) and is realisable, its separator cells are among $A$'s
cells when $A$ is small; in the boundary case the inequality
$\sigma(A)>\thick_\Part(A)$ is exactly the statement that interior content is
present. Finally, if $\sigma(A)\le\thick$ then since
$\thick\le\thick_\Part(A)$ would force $\sigma(A)\le\thick_\Part(A)$, the
interior content $\sigma(A)-\mu(A\cap\Sigma_\Part(A))$ cannot be positive, and
$A\subseteq\Sigma_\Part(A)$: $A$ is engulfed.
\end{proof}

\begin{remark}[Distinguishability is relative to the separator, not absolute]
\label{rem:relative}
The theorem locates distinguishability not in the absolute size of $A$ but in the
comparison $\sigma(A)$ versus $\thick_\Part(A)$. A region is silent---fails to
register---precisely when its separating boundary is thick relative to it. This
is the exact inversion the introduction promised: smallness is never itself the
cause of silence; being engulfed by one's own boundary is.
\end{remark}

\subsection{The sorites as a corollary}

We model an aggregate as a region grown by accretion of sub-floor grains and ask
when it becomes distinguishable. The classical sorites asks for the single grain
whose addition first makes the aggregate a heap. We show no such grain exists,
and locate the reason in the floor.

\begin{construction}[Accretive aggregate]
\label{con:accretion}
Let $A_0\subset A_1\subset A_2\subset\cdots$ be realisable regions with
$A_{n+1}=A_n\cup\cell_{i_n}$ formed by adjoining a single cell of measure
$\mu(\cell_{i_n})\le\thick$ \textup{(}a sub-floor grain\textup{)}. Write
$\tau_n:=\tau_\Part(A_n)$ for the relative thickness at stage $n$.
\end{construction}

\begin{corollary}[No sharp threshold; the sorites dissolves]
\label{cor:sorites}
For the accretive aggregate of \Cref{con:accretion}:
\begin{enumerate}[label=\textup{(\roman*)},leftmargin=2.2em]
\item \textbf{No single grain is decisive.} The transition from engulfed
\textup{(}$\tau_n\ge1$\textup{)} to distinguishable \textup{(}$\tau_n<1$\textup{)}
cannot be effected by any single sub-floor grain alone: each grain changes
$\sigma(A_n)$ by at most $\thick$ and $\thick_\Part(A_n)$ by a bounded amount, so
$\tau_n$ moves continuously and no one step carries it across $1$ from a value
bounded away from $1$.
\item \textbf{Distinguishability is emergence across a band.} There is a range of
stages over which $\tau_n$ descends through a neighbourhood of $1$; the aggregate
becomes distinguishable somewhere in this band, whose width is set by the floor
$\thick$, not at any point.
\item \textbf{The grain is silent because its boundary is thick relative to it.}
Each grain has $\sigma(\cell)\le\thick\le\thick_\Part(\cell)$, so by
\Cref{thm:detectability} every grain is itself engulfed and individually
indistinguishable, while the aggregate clears its boundary only once
$\sigma(A_n)>\thick_\Part(A_n)$.
\end{enumerate}
Hence the sorites has no decisive grain not because the predicate is vague in a
logical sense but because distinguishability is the crossing of a positive-width
band $\tau\approx1$, and a sub-floor grain cannot, by itself, carry an aggregate
across a band wider than its own contribution.
\end{corollary}

\begin{proof}
Each adjunction increases $\sigma$ by $\mu(\cell_{i_n})\le\thick$ and changes the
separator by at most the measure of the cells whose bordering status flips, a
quantity bounded in terms of $\thick$ by \Cref{thm:thickness}. Hence
$\tau_{n+1}-\tau_n$ is bounded in magnitude by a quantity controlled by
$\thick/\sigma(A_n)$, which is $<1$ once $\sigma(A_n)>\thick$. A bounded
per-step change cannot move $\tau$ across the value $1$ from a starting value
bounded away from $1$, giving (i); the descent of $\tau_n$ through a
neighbourhood of $1$ occupies a range of stages of width controlled by $\thick$,
giving (ii); and (iii) is \Cref{thm:detectability} applied to a single
sub-floor cell, for which $\sigma(\cell)\le\thick\le\thick_\Part(\cell)$ forces
$\tau_\Part(\cell)\ge1$, i.e.\ engulfment.
\end{proof}

\begin{remark}[Relation to theories of vagueness]
\label{rem:vagueness}
The standard responses to the sorites---degree-theoretic~\citep{zadeh1965},
supervaluationist~\citep{fine1975,keefe2000}, and
epistemicist~\citep{williamson1994}; see~\citet{black1937,hyderaffman2018}---all
locate the difficulty in the semantics of the predicate. The present account is
orthogonal and geometric: it places the difficulty in the object, namely in the
positive content of the separator that any distinguishing of the aggregate must
resolve. There is no hidden sharp threshold (against epistemicism), no
many-valued truth (against degree theories), and no need for admissible
precisifications (against supervaluation); there is a band of positive width
$\thick$ across which a region emerges from its boundary. The dissolution is a
theorem about separators, not a logic of partial truth.
\end{remark}

\begin{figure*}[!htbp]
    \centering
    \includegraphics[width=\textwidth]{figures/panel_4.png}
    \caption{\textbf{The sorites dissolved: emergence across a band, with no decisive grain.}
    (\textbf{A})~Relative thickness $\tau=\beta/\sigma$ on the linear $y$-axis ($0.03$ to $4$) versus the number of sub-floor grains added on the linear $x$-axis ($0$ to $142$), for an aggregate grown one unit at a time on a $12\times12$ instance. The curve starts well above the threshold $\tau=1$ (red)---a tiny aggregate is all boundary---and descends smoothly through it; there is no single grain at which $\tau$ leaps across $1$, the content of Corollary~\ref{cor:sorites}(i).
    (\textbf{B})~The same trajectory plotted as $\tau$ on the linear $y$-axis ($0.03$ to $4$) versus aggregate scale $\sigma(A_{n})$ on the linear $x$-axis ($1$ to $143$), threshold $\tau=1$ marked (red). Distinguishability is reached once $\sigma$ exceeds the boundary, across a band of positive width set by the floor, not at a point (Corollary~\ref{cor:sorites}(ii)).
    (\textbf{C})~Per-step change $|\Delta\tau|$ on the linear $y$-axis ($0$ to $1.33$) versus grains added on the linear $x$-axis ($0$ to $141$). The increments are bounded and decay once the aggregate clears its boundary, so no individual sub-floor grain carries the aggregate across the threshold: the grain is silent because its own boundary is thick relative to it (Corollary~\ref{cor:sorites}(iii)).
    (\textbf{D})~Three-dimensional overlay of three independent accretion orders (random orderings of the same atoms), each a curve of $\tau$ over aggregate scale $\sigma$ ($1$ to $\sim\!143$) at a fixed order-index $y$-coordinate. All three descend through the same $\tau\approx1$ band regardless of accretion order, showing the dissolution is a property of the geometry, not of any particular sequence of additions. Viewing angle: elevation $20^{\circ}$, azimuth $-62^{\circ}$.}
    \label{fig:panel-sorites}
\end{figure*}

% =====================================================================
\section{Locating versus visiting}
\label{sec:locating}
% =====================================================================

The results so far concern an agent that \emph{individuates} a region: draws a
separator, pays the floor, distinguishes the region from its complement. There is
a second, entirely different relation a region can stand in---being \emph{arrived
at} by a traversal of the space, with no separator drawn and no agent acting. This
section formalises the two relations and proves they obey opposite economies. The
distinction is purely measure- and order-theoretic; it names no dynamics. Its
consequence is a precise account of why a region can be free to be reached yet
costly to single out, and a classification of three invalid inferences that
silently exchange the two.

\subsection{Log-measure, probability, and traversal}

We add three abstract quantities, each with a standard home in measure theory and
information theory~\citep{shannon1948,cover2006}, and one abstract map.

\begin{definition}[Log-measure and probability of a region]
\label{def:logmeasure}
For a region $W\subseteq\Space$ with $\mu(W)>0$, its \emph{probability} is the
normalised measure
\[
p(W)\;:=\;\frac{\mu(W)}{\mu(\Space)}\in(0,1],
\]
and its \emph{log-measure} is
\[
S(W)\;:=\;\log\mu(W).
\]
A region is \emph{rare} if $p(W)$ is small, equivalently if $S(W)$ is small
relative to $\log\mu(\Space)$.
\end{definition}

\begin{remark}[One quantity under two names]
\label{rem:one-quantity}
$S(W)$ and $\log p(W)=S(W)-\log\mu(\Space)$ differ by an additive constant
independent of $W$. ``Small log-measure'' and ``small probability'' are therefore
not two correlated attributes of a region but a single attribute under two names.
No operation lowers $p(W)$ while holding $S(W)$ fixed, or raises one while fixing
the other; they move together by definition. This identity is used in
\Cref{prop:cannot-raise} and is the abstract counterpart of the familiar
log-count relation in information theory.
\end{remark}

\begin{definition}[Measure-preserving traversal]
\label{def:traversal}
A \emph{traversal} of $\Space$ is a measurable map $T\colon\Space\to\Space$ with
$\mu(T^{-1}E)=\mu(E)$ for every Borel $E$ \textup{(}it preserves $\mu$\textup{)}.
A point $x$ \emph{visits} a region $W$ \emph{at step} $n$ if $T^{n}(x)\in W$. The
\emph{visiting frequency} of $W$ along $x$ is
\[
\bar f_W(x)\;:=\;\lim_{N\to\infty}\frac1N\sum_{n=0}^{N-1}\mathbf 1_W\!\bigl(T^{n}(x)\bigr),
\]
when the limit exists. No agent, separator, or choice appears in this definition.
\end{definition}

\begin{remark}[Why this is not ``dynamics'']
\label{rem:not-dynamics}
A traversal is an iterated measure-preserving self-map---an order-theoretic and
measure-theoretic object, the discrete shift along the orbit
$x,Tx,T^2x,\dots$. We use it only through two classical facts: that on a
finite-measure space such a map returns almost every point to any positive-measure
region (recurrence), and that the visiting frequency, when it exists, equals the
region's probability (the ergodic average). Both are theorems of measure theory;
neither requires a continuum parameter, a rate, or any external scale.
\end{remark}

\subsection{Two relations to a region}

\begin{definition}[Visiting and locating]
\label{def:visit-locate}
A point \emph{visits} a region $W$ under a traversal $T$ when some iterate
$T^{n}(x)$ lies in $W$ \textup{(}\Cref{def:traversal}\textup{)}: a fact about the
orbit, involving no separator and no agent. An embedded agent \emph{locates} $W$
when it individuates $W$---draws a realisable separator marking $W$ off from
$\Space\setminus W$ to its finite resolution \textup{(}\Cref{def:individuation},
\Cref{def:realisable}\textup{)}: an act carrying the floor residue. Visiting
relates a traversal to a point; locating relates an agent to a region.
\end{definition}

\begin{remark}[The relations are categorically distinct]
\label{rem:two-relations}
At every step a point occupies exactly one maximally fine position, visited
because the traversal took it there, with no agent locating it. Conversely an
agent may locate a region no orbit visits within any finite horizon. A region's
being visited says nothing about its being located, and vice versa. The rest of
the section is the consequence of holding the two apart.
\end{remark}

\subsection{The asymmetry}

\begin{theorem}[Locating--Visiting Asymmetry]
\label{thm:asymmetry}
Let $W$ be a realisable region of $(\Space,d,\mu)$ under \Cref{ax:brs} with
$\Space$ connected, and let $T$ be a measure-preserving traversal. Then:
\begin{enumerate}[label=\textup{(\roman*)},leftmargin=2.4em]
\item \textbf{Locating is costly.} Any individuation of $W$ by an embedded agent
deposits separator content $\ge\thick>0$ \textup{(}\Cref{thm:thickness}\textup{)};
the cost does not depend on $p(W)$ and is never zero.
\item \textbf{Visiting is free.} Almost every point visits $W$ under $T$, with no
agent and no separator; by recurrence on the finite-measure space
$\mu(\Space)<\infty$ the visit occurs, and being a fact about the orbit it carries
no individuation cost.
\item \textbf{Only frequency scales with rarity.} The visiting frequency satisfies
$\bar f_W(x)=p(W)$ for an ergodic traversal \textup{(}and is bounded by the
recurrence fraction in general\textup{)}; thus a rare $W$ is visited just as
freely as a common one, only \emph{less often}. Rarity is a statement about
$p(W)$, never about the cost or possibility of a visit.
\end{enumerate}
\end{theorem}

\begin{proof}
(i) is \Cref{thm:thickness}: a realisable separator in a connected bounded
finitely-resolved space has measure $\ge\mu_{\min}=\thick>0$, independently of
$\mu(W)$. (ii): on a space with $\mu(\Space)<\infty$ a measure-preserving $T$
returns almost every point of any positive-measure region to that region (the sets
$T^{-k}W$ cannot be pairwise disjoint modulo null sets, since their equal measures
$\mu(W)>0$ would otherwise sum past $\mu(\Space)$); the visit is a property of the
orbit and involves no separator, hence no floor. (iii): the visiting frequency is
the average of the indicator $\mathbf 1_W$ along the orbit, which for an ergodic
$T$ equals $\int \mathbf 1_W\,\mathrm d(\mu/\mu(\Space))=p(W)$; smallness of the
frequency is exactly smallness of $p(W)$, and is independent of (i)--(ii).
\end{proof}

\begin{remark}[The asymmetry in one line]
\label{rem:needle}
A region may be hard to single out and yet effortlessly reached: it is hard to
\emph{locate} a distinguished region among many alternatives, but a traversal
passes through it regardless of whether anything locates it, and through a rare
region only \emph{less frequently}, never at greater cost. The cost lives entirely
on the locating side; the rarity lives entirely on the visiting side; nothing
carries one to the other.
\end{remark}

\subsection{The cost of locating among alternatives}

Locating a single region costs the floor (\Cref{thm:thickness}). Locating one
region \emph{among many} costs more, and the excess grows with how rare the target
is---a quantitative refinement of detectability.

\begin{theorem}[Locating cost grows with the log of the alternatives]
\label{thm:locating-cost}
Let $W$ be one region among $K$ pairwise-disjoint realisable alternatives of equal
measure that together exhaust a fixed portion of $\Space$. To single out $W$ from
the other $K-1$, an embedded agent performs at least $\lceil\log_2 K\rceil$ binary
individuations, each depositing residue $\ge\thick$; hence the total locating cost
is at least
\[
\thick\,\lceil\log_2 K\rceil .
\]
Since the target's probability is $p(W)=1/K$ of the exhausted portion, the cost
grows like $\thick\,\log_2(1/p(W))$: rarer targets are costlier to locate, while
\textup{(}by \Cref{thm:asymmetry}\textup{)} no rarer to visit.
\end{theorem}

\begin{proof}
Distinguishing $W$ from $K-1$ equally-sized alternatives is a search whose outcome
selects one of $K$ possibilities; any decision procedure resolving one of $K$
equiprobable outcomes makes at least $\lceil\log_2 K\rceil$ binary distinctions
(a counting bound: $b$ binary tests distinguish at most $2^{b}$ cases, so
$2^{b}\ge K$). Each binary distinction is an individuation---a separator between
the surviving candidates and the rejected ones---carrying residue $\ge\thick$ by
\Cref{thm:thickness}. Summing over the $\ge\lceil\log_2 K\rceil$ distinctions
gives total residue $\ge\thick\lceil\log_2 K\rceil$. With $p(W)=1/K$ this is
$\thick\,\lceil\log_2(1/p(W))\rceil$.
\end{proof}

\begin{remark}[Consistency with the cardinality derivation]
\label{rem:consistent-card}
\Cref{thm:locating-cost} refines \Cref{prop:card}: where the cardinality argument
showed that exhausting the regions of $\Space$ requires unbounded refinement, each
stage carrying its own floor, this shows that even \emph{within} a fixed finite
family of $K$ alternatives the cost of singling out one member scales with
$\log K$. The floor appears once per binary distinction, and the number of
distinctions is set by the rarity of the target.
\end{remark}

\subsection{Three invalid transfers}

The asymmetry blocks a family of inferences that move an attribute of locating
onto visiting, or silently alter the measure that constitutes a region. We state
the three that arise in practice and show each is invalid; they share a single
form (\Cref{prop:one-form}).

\begin{proposition}[Naming confers no structural priority]
\label{prop:rename}
Let $y=T^{n}(x)$ with $n\ge1$, so $y$ is reached by the traversal and has
predecessors. Designating $y$ ``the starting point'' is a locating act
\textup{(}a label\textup{)} and does not endow $y$ with the structural property
of an origin---having no predecessor. Any statement whose hypothesis is ``has no
predecessor'' fails to apply to $y$ whatever $y$ is called.
\end{proposition}

\begin{proof}
``Starting point'' has two senses: a bookkeeping origin chosen by the agent, and
the structural origin, a point with no predecessor under $T$. A reached point
$y=T^{n}(x)$, $n\ge1$, has the predecessor $T^{n-1}(x)$, so it meets the first
sense by stipulation and fails the second by construction. A statement proved
under ``no predecessor'' reads a property $y$ lacks; the label supplies only the
inert first sense. The transfer carries nothing.
\end{proof}

\begin{proposition}[Specificity does not entail an arranger]
\label{prop:arrange}
From ``$W$ is highly specific'' \textup{(}rare; small $p(W)$\textup{)} one cannot
infer ``some agent located or produced $W$.'' Specificity is a measure attribute
of the region; production is an act of an agent. The first does not entail the
second.
\end{proposition}

\begin{proof}
\Cref{thm:asymmetry}(ii) provides a standing counterexample: under any traversal,
the position $T^{n}(x)$ at each step is a maximally fine, individually rare
location that no agent located---it is visited, not located. Hence a region can be
as rare as one likes while having been singled out by no one. Rarity is $p(W)$;
arrangement is agency; rarity does not imply agency.
\end{proof}

\begin{proposition}[A region cannot be made probable at fixed log-measure]
\label{prop:cannot-raise}
There is no operation that makes a rare region the most probable destination while
it remains rare. ``Most probable'' is largest measure and ``rare'' is smallest
measure; by \Cref{rem:one-quantity} these are the same scale, so the demand is that
one region be simultaneously least and most. Any construction appearing to achieve
it has enlarged the region---changed $\mu(W)$, hence $S(W)$ and $p(W)$ together---and
so has replaced the rare region by a different, common one.
\end{proposition}

\begin{proof}
By \Cref{rem:one-quantity}, $p(W)$ and $S(W)$ are the same scale up to an additive
constant; maximising $p(W)$ maximises $\mu(W)$. To keep $W$ rare is to keep
$\mu(W)$ small; to make it most probable is to make $\mu(W)$ largest. Requiring
both is requiring $\mu(W)=\min$ and $\mu(W)=\max$, impossible unless the family is
trivial. A purported construction that ``raises the probability of the rare
region'' has increased $\mu(W)$ and thereby exchanged $W$ for a larger region; the
original rare region was not made probable but replaced.
\end{proof}

\begin{proposition}[The three transfers share one form]
\label{prop:one-form}
\Cref{prop:rename,prop:arrange,prop:cannot-raise} are instances of one invalid
move: an attribute established for the agent's relation to a region
\textup{(}locating: naming, singling out, producing\textup{)} is carried onto the
traversal's relation to the region \textup{(}visiting\textup{)}, or the measure
that constitutes the region is silently changed. Under \Cref{thm:asymmetry} and
\Cref{rem:one-quantity} each move is blocked, because the two relations obey
opposite economies and the constituting measure is the probability.
\end{proposition}

\begin{proof}
\Cref{prop:rename} carries a label (locating) onto a structural predicate of
visiting (having no predecessor). \Cref{prop:arrange} carries an agency claim
(production) onto a measure attribute (rarity). \Cref{prop:cannot-raise} alters
the measure that is the probability. In each case validity would require visiting
to inherit locating's cost-structure, or a region to change its measure while
keeping its identity; \Cref{thm:asymmetry} forbids the former and
\Cref{rem:one-quantity} the latter.
\end{proof}

\begin{remark}[What this section adds]
\label{rem:locating-adds}
The floor (\Cref{sec:thickness,sec:floor}) priced the act of individuation; this
section shows that price is paid by \emph{locating} alone and not by
\emph{visiting}, that it scales with the log of a target's rarity
(\Cref{thm:locating-cost}), and that the appearance of paradox in problems about
rare configurations is produced by transferring the locating price onto visiting
(\Cref{prop:one-form}). A reader will recognise that many concrete problems about
rare states---in which a configuration seems impossible to reach because it is
costly to arrange---have exactly this shape; the resolution in each case is the
asymmetry proved here, that what is costly to single out may be free to be reached
and merely seldom so. The self-referential obstruction of the next section shows
why even the locating price can never be driven to zero: the residue it leaves
cannot be certified away from within.
\end{remark}

\begin{figure*}[!htbp]
    \centering
    \includegraphics[width=\textwidth]{figures/panel_5.png}
    \caption{\textbf{Locating versus visiting: cost rises with rarity; being reached does not.}
    (\textbf{A})~Cost on the linear $y$-axis ($0$ to $8$) versus number of alternatives $K$ on a logarithmic $x$-axis ($2$ to $256$). The locating cost (blue step) follows $\beta\lceil\log_{2}K\rceil$ from Theorem~\ref{thm:locating-cost}, while the visiting cost (red) is identically zero for all $K$: singling a region out is priced and the price grows with the alternatives, whereas a traversal's passing through it is free (Theorem~\ref{thm:asymmetry}).
    (\textbf{B})~Locating cost on the linear $y$-axis ($1$ to $8$) versus $\log_{2}(1/p)$ on the linear $x$-axis ($1$ to $8$), where $p=1/K$ is the target's probability. The relation is linear with slope $\beta$: the cost of locating scales with the logarithm of the target's rarity, the quantitative refinement of detectability promised in Theorem~\ref{thm:locating-cost}.
    (\textbf{C})~Visiting frequency $\bar f_{W}$ on a logarithmic $y$-axis ($\sim\!3.9\times10^{-3}$ to $0.5$) versus region probability $p(W)$ on a logarithmic $x$-axis (same range). The data lie exactly on the identity $\bar f_{W}=p(W)$, the ergodic-average clause of Theorem~\ref{thm:asymmetry}(iii): rarity enters a visit only through its frequency, never through its cost or possibility.
    (\textbf{D})~Three-dimensional surface of locating cost $\beta\lceil\log_{2}K\rceil$ over $\log_{2}K$ ($1$ to $8$) and floor $\beta$ ($0.5$ to $2$), magma colourmap. The cost rises linearly in both $\log_{2}K$ and $\beta$ and is bounded away from zero whenever $\beta>0$: there is no free locating, and the dependence on rarity is logarithmic throughout. Viewing angle: elevation $22^{\circ}$, azimuth $-55^{\circ}$.}
    \label{fig:panel-locating-visiting}
\end{figure*}


% =====================================================================
\section{Non-self-verifiability and the irreducible residue}
\label{sec:residue}
% =====================================================================

The previous sections show the separator has positive content and that this
content silences sub-floor regions. We now reach the deepest obstruction: the
separator cannot be \emph{verified} from within. To verify a part is to decide it
against the totality that constitutes it---its negation-sequence---and that
decision is self-referential, with no consistent value. We prove this by a
self-contained diagonal argument (\Cref{thm:diagonal}), deduce that the
irreducible residue is not a locatable set but a self-referential impossibility
(\Cref{thm:residue-not-measure}), and conclude that the floor can only ever be
bounded below, never exhibited (\Cref{cor:floor-only-bounded}).

\subsection{Verification as deciding membership in the negation}

Recall (\Cref{def:negation}) that $A$ is constituted as $\Space\setminus\Neg(A)$:
its identity is fixed by the totality $\Neg(A)=\Space\setminus A$ of what it is
not. We make verification precise as the decision that constitutes this identity.

\begin{definition}[Verifier]
\label{def:verifier}
A \emph{verifier} for the parts of $\Space$ is a map $\Ver$ that, given a part
$X$ and a candidate constituting totality $T$, returns $\Ver(X,T)\in\{0,1\}$,
intended to read $1$ iff $X$ is exactly the part constituted by withholding $T$,
i.e.\ iff $X=\Space\setminus T$. To \emph{verify $A$} is to evaluate
$\Ver(A,\Neg(A))$: to decide whether $A$ is the part its own negation-sequence
constitutes.
\end{definition}

\begin{remark}[Why verification is membership in the negation]
\label{rem:ver-membership}
By \Cref{def:negation}, $A=\Space\setminus\Neg(A)$ exactly when $A$ and $\Neg(A)$
partition $\Space$; deciding this is deciding, of the constituting totality
$\Neg(A)$, whether what it omits is precisely $A$---a question about $A$ posed
entirely in terms of the collection $\Neg(A)$ that is defined as
``everything but $A$.'' Verification of $A$ is thus inseparable from a question
of self-membership: $A$ figures, through its own definition, inside the totality
against which it is checked.
\end{remark}

\subsection{A self-contained diagonal obstruction}

We now prove, without invoking any external impossibility theorem, that no
verifier decides every part correctly. The argument is a diagonal one, in the
shape of the classical results of~\citet{russell1903,cantor1891,tarski1936} and
the general form of~\citet{lawvere1969}, but it is self-contained: it uses only
the definitions above.

Two ingredients are needed before the diagonal can be run, and we isolate them as
an explicit hypothesis, because whether they hold is exactly the delicate point.
The diagonal collection
\[
D\;:=\;\bigl\{\,X : \Ver(X,\Neg(X))=0\,\bigr\}
\]
is, a priori, only a \emph{collection of parts indexed by the verifier's
behaviour}; for the verifier to be asked about $D$ in the sense of the theory, $D$
must itself be an admissible object of the same kind---a region of $\Space$ on
which $\Ver$ is defined.

\begin{definition}[Self-applicable verifier]
\label{def:self-applicable}
A verifier $\Ver$ is \emph{self-applicable} if
\begin{enumerate}[label=\textup{(SA\arabic*)},leftmargin=3em]
\item\label{sa:total} it is total on the regions of $\Space$ (it returns a value
in $\{0,1\}$ on every region and its negation), and
\item\label{sa:closed} the diagonal collection $D$ above is itself a region of
$\Space$---a Borel set on which $\Ver(\,\cdot\,,\Neg(\,\cdot\,))$ is defined---so
that $\Ver(D,\Neg(D))$ is a legal query.
\end{enumerate}
\end{definition}

\begin{theorem}[Diagonal obstruction to verification]
\label{thm:diagonal}
No self-applicable verifier exists. Equivalently, there is no verifier $\Ver$
that is total on regions \ref{sa:total} and closed under its own diagonal
\ref{sa:closed}: for any such $\Ver$ the value $\Ver(D,\Neg(D))$ cannot be
consistently assigned, since $\Ver(D,\Neg(D))=1$ and $\Ver(D,\Neg(D))=0$ each
entail the other. In particular no verifier satisfying \ref{sa:closed} returns
$1$ on every region.
\end{theorem}

\begin{proof}
Suppose, for contradiction, a self-applicable verifier $\Ver$ exists. By
\ref{sa:closed} the collection $D$ is a region and $\Ver(D,\Neg(D))$ is a legal
query; by \ref{sa:total} it has a value. $D$ collects exactly those parts $X$ for
which the verifier returns $0$ on the self-referential query $\Ver(X,\Neg(X))$.
Now ask the verifier about $D$ itself.

If $\Ver(D,\Neg(D))=1$, then $D$ is declared constituted by its own negation;
but by the defining property of $D$, membership-defining condition for $D$ is
$\Ver(X,\Neg(X))=0$, so $D$ satisfying its own membership condition would require
$\Ver(D,\Neg(D))=0$. Reading $D$ through its own definition therefore forces
$\Ver(D,\Neg(D))=0$, contradicting the assumed value $1$.

If instead $\Ver(D,\Neg(D))=0$, then $D$ meets its own membership condition, so
$D$ is among the parts collected by $D$, i.e.\ $D$ is constituted as the part its
negation withholds, which is the assertion $\Ver(D,\Neg(D))=1$, contradicting the
assumed value $0$.

Both values are self-refuting, so no consistent value exists for
$\Ver(D,\Neg(D))$. This contradicts \ref{sa:total}, which requires a value. Hence
no verifier satisfies both \ref{sa:total} and \ref{sa:closed}; in particular none
satisfying \ref{sa:closed} returns $1$ on every region. This is the diagonal
obstruction.
\end{proof}

\begin{remark}[On the closure hypothesis \ref{sa:closed}]
\label{rem:self-app}
The theorem is exactly as strong as \ref{sa:closed}, and we are careful not to
claim more than we prove. \ref{sa:closed} asserts that the behaviour-defined
collection $D$ is again a region of $\Space$ on which the verifier is defined.
This is the genuine content of ``self-application'': the verifier can be asked
about an object built from its own verdicts. Two remarks on its status.

\emph{Why embedding makes it natural.} An agent that verifies regions of $\Space$
while being itself a part of $\Space$ is a plausible setting in which a
description of the agent's own verifying behaviour is realised as a region---the
agent's verdicts leave a trace in the same space they range over, and a region
selected by those verdicts is a region of that space. In such a setting
\ref{sa:closed} holds and the obstruction bites.

\emph{Why it is not automatic.} Nothing in \Cref{ax:brs} forces a
behaviour-defined collection to be Borel-encodable as a region. The classical
diagonal results show the shape of the argument but not that its preconditions are
met here: Russell's construction lives in unrestricted comprehension, which the
measurable setting does not grant; and the general fixed-point
theorem of~\citet{lawvere1969} delivers the contradiction only given a
point-surjection in a cartesian-closed category---structure that an embedding
\emph{may} supply but that must be exhibited, not assumed. We therefore treat
\ref{sa:closed} as a hypothesis, true in any model where the agent's verifying
behaviour is itself realisable as a region, and we mark every downstream claim
that depends on it. Discharging \ref{sa:closed} in full generality---constructing
the measurable encoding of parts on which $\Ver$ acts as a point-surjection---is a
separate undertaking we do not claim to have carried out.
\end{remark}

\subsection{The residue is not a set under the closure hypothesis}

The diagonal obstruction lets us characterise the irreducible residue, under the
same closure hypothesis \ref{sa:closed} that the obstruction requires. One might
hope that the un-verified remainder of an individuation is a definite set---small,
locatable, perhaps measurable---which a finer effort could in principle exhibit
and remove. We show that this hope is incompatible with \ref{sa:closed}: a
locatable residue would let one assemble a total verifier closed under its own
diagonal, which \Cref{thm:diagonal} forbids.

\begin{definition}[Irreducible residue]
\label{def:irreducible}
The \emph{irreducible residue} of the individuation of $A$ is the obstruction to
deciding $\Ver(A,\Neg(A))$: the content that an individuation cannot certify
about itself. It is \emph{exhibitable} if there is a realisable region $E$ with
$\mu(E)>0$ such that verifying \emph{every} region reduces to verifying it on
$\Space\setminus E$ together with a separate, always-terminating certification of
the residue localised in $E$.
\end{definition}

\begin{theorem}[The irreducible residue is not exhibitable, under \ref{sa:closed}]
\label{thm:residue-not-measure}
Let \Cref{ax:brs} hold and suppose the closure hypothesis \ref{sa:closed}---that
a behaviour-defined collection of regions is again a region. Then the irreducible
residue is not exhibitable: there is no realisable region $E$ of positive measure
whose separate certification completes the verification of every region. Under
\ref{sa:closed} the residue is a self-referential impossibility, not a locatable
set.
\end{theorem}

\begin{proof}
Suppose the residue were exhibitable by some realisable $E$ with $\mu(E)>0$:
certifying the residue localised in $E$ always terminates and, combined with
verification on $\Space\setminus E$, completes the verification of every region.
This combination is a verifier $\Ver^\ast$ that is total on regions
\ref{sa:total}. By the standing hypothesis \ref{sa:closed}, the diagonal
collection $D$ of $\Ver^\ast$ is again a region; so $\Ver^\ast$ satisfies both
\ref{sa:total} and \ref{sa:closed}, i.e.\ is self-applicable. \Cref{thm:diagonal}
shows no such verifier exists---a contradiction. Hence no exhibiting $E$ exists:
under \ref{sa:closed} the residue cannot be localised to a positive-measure set
and certified away. It is the standing impossibility of the self-referential
decision, not a region awaiting finer treatment.
\end{proof}

\begin{remark}[What survives if \ref{sa:closed} fails]
\label{rem:if-sa-fails}
If the closure hypothesis fails---if the agent's verifying behaviour is not
realisable as a region---then \Cref{thm:diagonal} and \Cref{thm:residue-not-measure}
do not apply, and the present section establishes nothing about exhibitability. The
paper's quantitative core is untouched by this: the floor $\thick>0$ is already
established three independent ways in \Cref{sec:thickness,sec:floor}, before any
diagonal argument is invoked, and the detectability, locating, and non-return
results of \Cref{sec:detect,sec:locating,sec:inquiry} rest on those derivations
alone. What rides on \ref{sa:closed} is solely the \emph{upgrade} from ``the floor
is positive'' to ``the residue is not a locatable set, only boundable below.'' We
state that upgrade as conditional and claim no more.
\end{remark}

\begin{corollary}[The floor can only be bounded below, under \ref{sa:closed}]
\label{cor:floor-only-bounded}
The resolution floor $\thick$ admits a strict positive lower bound
\textup{(}$\thick\ge\mu_{\min}>0$, \Cref{thm:thickness}\textup{)}
unconditionally. If in addition the closure hypothesis \ref{sa:closed} holds,
then $\thick$ has no exact exhibition as $\mu(R)$ for a definite residue set $R$:
every attempt to write the residue as the measure of a located set $R$ supplies
an exhibitable $E=R$, contradicting \Cref{thm:residue-not-measure}. In that case
the floor is known only as the positive quantity below which no individuation can
certify itself.
\end{corollary}

\begin{proof}
The lower bound is \Cref{thm:thickness}. If $\thick=\mu(R)$ for a realisable $R$
of positive measure, then $R$ exhibits the residue, contradicting
\Cref{thm:residue-not-measure}. Hence no such exhibition exists, and $\thick$ is
characterised solely as a positive lower bound on un-certifiable content.
\end{proof}

\begin{remark}[The recognition gap]
\label{rem:recognition-gap}
\Cref{thm:residue-not-measure}, where it applies, has a reading in terms of the
relation between an agent and the part it would verify. To certify $A$ completely, the agent must
decide $A$ against the totality $\Neg(A)$ that constitutes $A$; but the agent,
being embedded, is itself within that totality, so a complete certification would
require it to occupy both sides of the constituting partition at once---to be at
once the part verified and a portion of the totality against which it is
verified. The residue is the irreducible gap that keeps the two apart. Closing
the gap would collapse the distinction on which verification depends; the floor
$\thick$ is the measure of how far the gap can never be closed. This is the
precise sense in which the boundary of knowledge is constitutive, not contingent:
the un-knowable remainder is the very condition under which anything is
individuated at all.
\end{remark}


% =====================================================================
\section{Non-return and the cessation of inquiry}
\label{sec:inquiry}
% =====================================================================

We close the development with two order-theoretic results. The first
(\Cref{thm:nonreturn}) shows that the record of committed individuations is
monotone under every traversal, so no process returns to a prior individuation
state and resemblance never amounts to identity. The second
(\Cref{thm:cessation}) shows that, because thinning a separator has diverging
cost and bounded return, a finite agent halts at the floor; knowledge thereafter
attaches to bounded wholes and never to sub-floor parts. Neither result uses any
dynamics: both are statements about the order on committed separations.

\subsection{The committed record}

\begin{definition}[Committed record]
\label{def:record}
Fix a refinement chain of realisable partitions
$\Part_0\succeq\Part_1\succeq\cdots$. The \emph{committed record} after a finite
sequence of individuation steps is the count
\[
M\;:=\;\#\{\text{separator cells committed across all steps so far}\}\in\N,
\]
the number of cells moved out of an approximant's separator over the process
\textup{(}\Cref{def:process}, \Cref{cor:monotone-commit}\textup{)}.
\end{definition}

\begin{theorem}[Non-return]
\label{thm:nonreturn}
The committed record $M$ is monotone non-decreasing under every traversal of the
refinement chain. No traversal returns the record to a previously held smaller
value; in particular, a step intended to undo a prior commitment is itself a
committed step and increases $M$. Consequently two individuation states with the
same configuration of cells but different records are distinct, and resemblance
of configuration never entails identity of state.
\end{theorem}

\begin{proof}
By \Cref{cor:monotone-commit} each committing step increases the count of
committed separator cells by at least one, and the refinement order
$\Part_j\succeq\Part_{j+1}$ is one-directional: to ``undo'' a commitment is to
perform a further act---a re-coarsening or re-drawing---which by
\Cref{def:process} is itself a committing step and so increments $M$. Hence $M$
never decreases; any attempted return raises it. If two states $s,s'$ share a
configuration of cells but have records $M(s)\neq M(s')$, then they are
distinguished by the record itself: there is a partition of states, namely
``record $\le M(s)$ versus record $>M(s)$,'' separating them. By the principle
that states are identical only when no realisable partition distinguishes them,
$s\neq s'$. Resemblance of configuration is therefore compatible with difference
of record, and identity requires both.
\end{proof}

\begin{remark}[Resemblance versus identity]
\label{rem:resemblance}
\Cref{thm:nonreturn} is the order-theoretic core of irreversibility for
individuation: a process may revisit a configuration that \emph{looks like} an
earlier one, but the committed record that accompanies it has strictly grown, so
the state is new. The monotone record is the invariant that configuration alone
cannot capture. Nothing here invokes any flow or external parameter; the
monotonicity is a property of the order on commitments.
\end{remark}

\subsection{Cessation}

\begin{theorem}[Inquiry cessation]
\label{thm:cessation}
Let $\Gamma$ be a refinement cost \textup{(}\Cref{def:cost}\textup{)} and let the
return of an individuation---the additional content it correctly certifies---be
bounded by $\mu(\Space)<\infty$ \textup{(}\ref{brs:bdd}\textup{)}. Then for any
agent with a positive threshold on cost-per-return there is a finite thickness
$\thick_\ast>0$ at which further thinning is rejected: the agent halts at a
strictly positive floor and certifies nothing below it. Knowledge so obtained
attaches to bounded wholes of thickness $\ge\thick_\ast$ and never to sub-floor
parts.
\end{theorem}

\begin{proof}
By \Cref{thm:cost}, realising thickness $t$ costs $\Gamma_t\ge g(t)$ with
$g(t)\to+\infty$ as $t\to0^+$. The return is bounded above by $\mu(\Space)$, a
finite constant. Hence the marginal return per unit cost,
$\mu(\Space)/g(t)$, tends to $0$ as $t\to0^+$. An agent that continues thinning
only while cost-per-return stays below a fixed positive threshold $\theta$ stops
at the largest $t$ with $g(t)\le\mu(\Space)/\theta$; call it $\thick_\ast$. Since
$g$ is finite away from $0$, $\thick_\ast>0$. Below $\thick_\ast$ the agent does
not refine, so it certifies no region of thickness $<\thick_\ast$; by
\Cref{thm:detectability} regions of scale below the floor are in any case
engulfed. Thus knowledge attaches to wholes of thickness $\ge\thick_\ast$ and to
no sub-floor part.
\end{proof}

\begin{remark}[Why knowledge attaches to wholes]
\label{rem:wholes}
\Cref{thm:cessation} gives the final shape of the picture. Individuation is by
negation (\Cref{thm:negation-unique}); it is non-instantaneous and leaves a
residue (\Cref{thm:noninstant}); the residue cannot be exhibited or certified
(\Cref{thm:residue-not-measure}); the record of commitments only grows
(\Cref{thm:nonreturn}); and the cost of pursuing the residue to zero diverges
against a bounded return, so a finite agent rationally halts at a positive floor
(\Cref{thm:cessation}). What such an agent can know are bounded wholes---regions
that clear their own boundary thickness---attributed and reasoned about as units.
It cannot know its sub-floor parts, not for want of effort but because below the
floor a part is indistinguishable from its boundary and its certification is
self-referentially obstructed. The boundary of knowledge is the floor, and the
floor is positive.
\end{remark}

\begin{figure*}[!htbp]
    \centering
    \includegraphics[width=\textwidth]{figures/panel_6.png}
    \caption{\textbf{Non-return and cessation: the record only grows, and inquiry rationally halts.}
    (\textbf{A})~Committed record $M$ on the linear $y$-axis ($4$ to $378$) versus process step on the linear $x$-axis ($0$ to $60$). The record is monotone non-decreasing along the whole separation process, and the final marked point is an attempted ``undo'' that nonetheless increments $M$---the content of Theorem~\ref{thm:nonreturn}: a step intended to reverse a commitment is itself a committed step, so no traversal returns the record to a prior value.
    (\textbf{B})~The cost $g(t)=-\log t$ on the linear $y$-axis ($\sim\!0.12$ to $12.6$) versus separator thickness $t$ on a logarithmic $x$-axis ($\sim\!3.5\times10^{-6}$ to $0.89$), reproducing the divergence of Theorem~\ref{thm:cost} on the instance used for cessation.
    (\textbf{C})~Marginal return per unit cost on the linear $y$-axis ($\sim\!8$ to $869$) versus thickness $t$ on a logarithmic $x$-axis (same range), with several halting thicknesses $\beta_{\ast}=e^{-\mu(\Omega)/\theta}$ marked (red) for thresholds $\theta\in\{0.5,1,2,4\}$. Because the return is bounded by $\mu(\Omega)<\infty$ while the cost diverges, the marginal return falls toward zero and a finite agent with any positive cost-per-return threshold halts at a strictly positive $\beta_{\ast}$ (Theorem~\ref{thm:cessation}).
    (\textbf{D})~Three-dimensional surface of marginal return per cost over $\log_{10}t$ ($\sim\!-5.4$ to $-0.05$) and threshold $\theta$ ($0.5$ to $4$), viridis colourmap. The surface falls to a low plateau as $t\to0$ for every $\theta$: the diverging-cost, bounded-return economy forces a halt across the whole range of thresholds, no agent driving the separator to zero. Viewing angle: elevation $24^{\circ}$, azimuth $-58^{\circ}$.}
    \label{fig:panel-nonreturn-cessation}
\end{figure*}


% =====================================================================
\section{Numerical verification}
\label{sec:numerical}
% =====================================================================

The theorems are proved from the definitions and require no computation. It is
nonetheless instructive to exhibit them on concrete instances of the hypothesis,
both as a check against error and to display the quantitative shape of each
result. This section reports an exhaustive numerical witness: every theorem-group
is tested on a large family of explicitly constructed bounded resolvable spaces,
and the measured quantities are plotted against the bounds the theorems predict.

\subsection{Method}

A bounded resolvable space is realised as a finite metric--measure space: a finite
integer-lattice box, a counting measure with a fixed atom weight, and a partition
into axis-aligned cells of side $\ge\reso$. By construction such an instance
satisfies \Cref{ax:brs}---it is bounded with finite measure, every cell has
diameter at least the resolution, and finer boxes refine it---and an embedded
agent on it can name only unions of cells, exactly the operational restriction of
\Cref{def:realisable}. On each instance the realisable regions, their separators,
boundary thickness, scale, log-measure, and the committed record are computed
directly from the definitions, and each theorem's inequality or equivalence is
checked.

A randomised sweep drew instances of varying dimension, extent, cell size, and
atom weight, and ran all seven theorem-group checks on each. Across
$500$ instances the suite performed $96{,}024$ individual checks, all of which
passed: no realisable separator was sharp; every boundary thickness met
$\thick\ge\mu_{\min}$; the three floor derivations agreed on every instance;
distinguishability matched $\sigma>\thick_A$ in every region tested; the
self-referential value admitted no consistent assignment; and the committed
record was monotone on every traversal. The figures below are drawn from these
same computations.


% =====================================================================
\section{Scope and conclusion}
\label{sec:scope}
% =====================================================================

\subsection{What has been proved}

From the single hypothesis of a bounded resolvable space (\Cref{ax:brs}) we have
established the following, each from stated definitions:

\begin{enumerate}[label=(\arabic*),leftmargin=2.4em]
\item \textbf{Individuation by negation} (\Cref{thm:negation-unique}): in a whole
with no distinguished element, a part can be fixed only as the complement of its
negation-sequence; this is the unique selector-free individuation.
\item \textbf{Non-instantaneity and residue} (\Cref{thm:noninstant}): no
individuation of a genuine part completes in zero steps, and at every stage before
completion a separator of measure $\ge\thick$ remains uncommitted.
\item \textbf{Boundary thickness} (\Cref{thm:thickness}): every realisable
separator has measure $\ge\mu_{\min}>0$; no sharp cut exists for an embedded
agent. The floor $\thick$ is strictly positive.
\item \textbf{Three agreeing derivations} (\Cref{thm:floor-agree}): the
geometric, representational, and cost-limited floors are all positive and bound
the same constant.
\item \textbf{Detectability and the sorites} (\Cref{thm:detectability},
\Cref{cor:sorites}): a region is distinguishable iff its scale exceeds its
boundary thickness; the sorites dissolves as emergence across a band of width
$\thick$, with no decisive grain.
\item \textbf{Locating versus visiting} (\Cref{thm:asymmetry},
\Cref{thm:locating-cost}): individuating a region is costly and the cost grows
with the logarithm of the region's rarity, while a measure-preserving traversal
visits the region freely and scales with rarity only in frequency; three invalid
inferences transferring the former onto the latter are classified and blocked
(\Cref{prop:one-form}).
\item \textbf{Non-self-verifiability} (\Cref{thm:diagonal},
\Cref{thm:residue-not-measure}): no verifier is total and closed under its own
diagonal, an unconditional self-contained result; and \emph{under the explicit
closure hypothesis} \ref{sa:closed}---that a behaviour-defined collection of
regions is again a region---the irreducible residue is not an exhibitable set and
the floor can only be bounded below, never displayed. This last upgrade is stated
conditionally; the positivity of the floor (items 3--4) does not rely on it.
\item \textbf{Non-return and cessation} (\Cref{thm:nonreturn},
\Cref{thm:cessation}): the committed record is monotone, so no process returns to
a prior individuation state; thinning a separator has diverging cost and bounded
return, so a finite agent halts at the floor and knows only bounded wholes.
\end{enumerate}

These results turn on one constant, the floor $\thick>0$. Its positivity is
reached from four independent directions---topological, representational,
cardinality-theoretic, and economic---each shown to bound or price the same
quantity unconditionally. Two further directions speak to its character rather
than its positivity: the self-referential, which under the closure hypothesis
\ref{sa:closed} shows the floor cannot be displayed as a located set, and the
order-theoretic, which shows the record of individuations cannot return.

\subsection{Scope}

The theorems are statements about an abstract bounded resolvable space and an
embedded agent that individuates parts within it by negation. We have made no
claim about any particular system, and the proofs use only the metric--measure
and order structure declared in \Cref{sec:setting}. The hypotheses
(\Cref{ax:brs}) are deliberately minimal: boundedness with compact closure,
finite resolution, and hierarchical refinability. Where each clause is used has
been marked, and the necessity of boundedness is shown by a counterexample
(\Cref{prop:bdd-necessary}). Material that could not be proved at this standard
has been omitted rather than asserted.

One hypothesis beyond \Cref{ax:brs} is used, and only in \Cref{sec:residue}: the
closure hypothesis \ref{sa:closed}, that a collection of regions defined by a
verifier's own behaviour is again a region of $\Space$. We do not prove
\ref{sa:closed}; we isolate it, argue that an embedding makes it natural, note
that discharging it in general would require exhibiting the measurable encoding
and point-surjection that the classical fixed-point theorem of~\citet{lawvere1969}
needs, and mark every conclusion that depends on it. The positive floor and all
results of \Cref{sec:thickness,sec:floor,sec:detect,sec:locating,sec:inquiry}
stand independently of \ref{sa:closed}.

\subsection{Conclusion}

A bounded whole, finitely resolved, cannot be divided into a part and its
complement without depositing positive content into the separator, and that
content cannot be driven to zero---and, where the agent's own verifying behaviour
is itself a region, cannot be certified from within. The part is
individuated by what it is not; the separation is never instantaneous and never
exact; the residue is irreducible; the record of separations only grows; and a
finite agent rationally halts at a positive floor, knowing bounded wholes and not
their sub-floor parts. The floor $\thick>0$ is the single invariant through which
individuation, boundary thickness, detectability, verification, and non-return
all speak. It is the irreducible boundary of a bounded resolvable space.

% =====================================================================
\bibliographystyle{plainnat}
\bibliography{references}
% =====================================================================

\end{document}
